\documentclass[11pt,oneside]{report}
\usepackage[letterpaper,margin=0.86in,headheight=15pt]{geometry}
\usepackage[T1]{fontenc}
\usepackage{lmodern,microtype}
\usepackage{amsmath,amssymb,mathtools,bm}
\usepackage{graphicx,booktabs,array,longtable}
\usepackage{xcolor}
\usepackage[font=small,labelfont=bf]{caption}
\usepackage{fancyhdr}
\usepackage{enumitem}
\usepackage{listings}
\usepackage{hyperref}
\usepackage[nameinlink,noabbrev]{cleveref}
\definecolor{accent}{HTML}{1F6F5F}
\hypersetup{colorlinks=true,linkcolor=accent,citecolor=accent,urlcolor=accent,
pdftitle={From Bloch Bands to Floquet Bands: Controlling Atomic Transport in a One-Dimensional Optical Lattice},pdfauthor={Lakshya Bhardwaj}}
\setlength{\parskip}{4pt}
\setlength{\parindent}{1.2em}
\setlength{\emergencystretch}{1em}
\newcolumntype{P}[1]{>{\raggedright\arraybackslash}p{#1}}
\setlist{itemsep=3pt,topsep=4pt}
\pagestyle{fancy}\fancyhf{}\fancyhead[L]{\small From Bloch bands to Floquet bands}\fancyhead[R]{\small\thepage}
\renewcommand{\headrulewidth}{0.3pt}
\fancypagestyle{plain}{\fancyhf{}\fancyfoot[C]{\thepage}\renewcommand{\headrulewidth}{0pt}}
\numberwithin{equation}{chapter}
\graphicspath{{figures/}}
\newcommand{\ER}{E_{\mathrm R}}
\newcommand{\Jeff}{J_{\mathrm{eff}}}
\newcommand{\ii}{\mathrm i}
\newcommand{\dd}{\mathrm d}
\newcommand{\avg}[1]{\left\langle #1\right\rangle}
\newcommand{\ket}[1]{\lvert #1\rangle}
\newcommand{\bra}[1]{\langle #1\rvert}
\newcommand{\braket}[2]{\langle #1\vert #2\rangle}
\newcommand{\fig}[3]{\begin{figure}[!htbp]\centering\includegraphics[width=\linewidth]{#1.pdf}\caption{#2}\label{#3}\end{figure}}
\lstset{basicstyle=\ttfamily\small,breaklines=true,frame=single,rulecolor=\color{black!20},columns=fullflexible}
\begin{document}
\begin{titlepage}
\centering
\vspace*{1.0in}
{\LARGE\bfseries From Bloch Bands to Floquet Bands:\par}
\vspace{0.2in}
{\Large Controlling Atomic Transport in a\par One-Dimensional Optical Lattice\par}
\vspace{0.5in}
{\large Lakshya Bhardwaj\par}
{Department of Physics, Arizona State University\par}
\vspace{0.2in}
{Course project in quantum theory of solids\par}
\vspace{0.6in}
\begin{minipage}{0.87\textwidth}
\textbf{Abstract.} A spatially periodic optical potential produces Bloch bands whose slopes govern atomic motion. Periodic lattice shaking replaces these energy dispersions by quasienergy dispersions. Starting from the optical cosine potential, this report derives the stationary plane-wave eigenproblem, the transformations to an inertial-force description, and the Floquet--Bloch equations. An exactly solvable nearest-neighbor model gives $\Jeff=J\mathcal J_0(K)$, where $K$ is the force amplitude times the lattice spacing divided by the driving quantum. This relation explains band narrowing, complete band collapse within that model, and inversion of the dispersion. Exact-model wave-packet propagation distinguishes reversible motion within each cycle from localization at complete cycles, and a packet with fixed nonzero central quasimomentum reveals the sign of $\Jeff$. Reproduced continuum bands quantify longer-range hopping corrections, while a bounded multiband calculation demonstrates resonant failure of the isolated-band description. The results are connected to the measurements of Lignier et al. (2007) and Eckardt et al. (2009). All figures are generated by the accompanying commented Python files, with numerical benchmarks and convergence checks.
\end{minipage}
\vfill
{\small Foundation: Martin Holthaus, J. Phys. B \textbf{49}, 013001 (2016).\par}
{\small Model calculations; no experimental data are fitted or presented as new measurements.\par}
\end{titlepage}
\pagenumbering{roman}
\tableofcontents
\clearpage
\chapter*{Scope and conventions}
\addcontentsline{toc}{chapter}{Scope and conventions}
The report is self-contained at the level of single-particle quantum mechanics, Fourier analysis and elementary solid-state band theory. Its principal foundation is Holthaus's tutorial on quasienergy bands \cite{Holthaus2016}. The tight-binding identities, dynamic-localization mechanism and experimental phenomena are established results. Their derivations and the numerical executions here are reproductions for a course project. The selected comparisons involving longer-range hopping and resonant excitation are additional teaching calculations, with no claim of novelty.

Three levels of description are kept distinct throughout:
\begin{enumerate}
\item The \emph{continuous optical potential} includes several bands and is truncated only for computation in a plane-wave basis.
\item The \emph{isolated lowest band} retains its complete, generally noncosine dispersion but omits transitions to other bands.
\item The \emph{nearest-neighbor model} additionally discards all hopping matrix elements except those between adjacent sites. Its driven dynamics can be solved exactly without a high-frequency expansion.
\end{enumerate}
The main calculations neglect atom--atom interactions, longitudinal confinement, disorder, dissipation and laser noise. A one-dimensional Hamiltonian describes motion along the lattice; independent transverse confinement is assumed to prevent transverse excitations. A dilute condensate can sometimes approximate the single-particle transport, but that approximation needs experimental justification.

The potential convention is
\begin{equation}\label{eq:convention}
 V(x)=V_0\sin^2(k_{\mathrm L}x)=\frac{V_0}{2}-\frac{V_0}{2}\cos(2k_{\mathrm L}x),\qquad V_0=s\ER>0.
\end{equation}
Holthaus uses $+(V_0/2)\cos(2k_{\mathrm L}x)$. Translating the origin by half a lattice spacing changes its cosine sign; adding $V_0/2$ gives \cref{eq:convention}. Energy gaps, bandwidths and transport are unaffected. Absolute energies in this report include the $V_0/2$ offset.

We use the lattice spacing $d=\pi/k_{\mathrm L}$, wave number $k$, dimensionless quasimomentum $q=kd$, recoil energy $\ER=\hbar^2 k_{\mathrm L}^2/(2m)$ and drive angular frequency $\omega$. The dimensionless drive strength is always $K=F_0d/(\hbar\omega)$; it must not be confused with an energy-valued forcing parameter also called $K$ in some references. Ordinary Bessel functions are denoted by $\mathcal J_n$ and hopping by $J$.
All transport centers and widths are measured relative to the lattice. The symbol $x$ in their moment notation abbreviates the moving coordinate $y$ used in the frame transformation; laboratory centers are explicitly labeled $x_{\rm lab}$.

\clearpage\pagenumbering{arabic}
\chapter{Introduction and motivation}
\section{An optical crystal for neutral atoms}
An optical lattice is a spatially periodic light-shift potential. Two counterpropagating laser beams form a standing wave. For an atom far from an internal optical resonance, the rapidly oscillating electromagnetic field can be eliminated, leaving a conservative potential proportional to its local intensity. Atoms then move through a periodic array of potential wells rather than through an array of material ions.

This construction realizes the central ingredient of electron band theory: a periodic potential acting on a quantum wave. Bloch states, reciprocal lattice vectors, Brillouin zones and energy gaps follow from the same spatial translation symmetry. The atom is neutral, however, and its mass, interactions, trapping environment and preparation differ from those of an electron in a solid. Optical lattices therefore simulate selected Hamiltonians rather than reproducing every microscopic aspect of a material. Laser intensity controls the lattice depth; relative optical phase controls its position; beam geometry sets its dimensionality. Interactions and particle statistics provide further control \cite{Jaksch1998,Bloch2008}.

Transport is a particularly clear observable. A localized packet occupies many quasimomenta. Different band energies give these components different dynamical phases, causing the packet to spread. A packet narrow in quasimomentum moves at the group velocity determined by the local band slope. A method that changes dispersion without changing the static lattice depth can therefore change spreading and directed motion.

\section{Why periodic driving changes the question}
For a stationary Hamiltonian, energy is conserved and the relevant spectral functions are $E_n(k)$. For a Hamiltonian with period $T=2\pi/\omega$, time translation by $T$ is the symmetry. Its eigenphases define quasienergies $\varepsilon_\alpha(k)$, which are equivalent modulo $\hbar\omega$. Floquet engineering uses controlled periodic modulation to shape these quasienergies and their states. The drive need not be a small perturbation: the ratio $F_0d/(\hbar\omega)$ can be of order unity even when the effective hopping is greatly reduced.

The project asks four central questions:
\begin{enumerate}
\item How does shaking change the width, slope and curvature of a Bloch band?
\item When does a driven wave packet stop spreading, and does it remain motionless within a driving cycle?
\item What distinguishes a narrowed band from an inverted band, and which observable reveals the sign of effective tunneling?
\item Which assumptions make the Bessel-function description valid, and how do longer-range hopping and interband transitions limit it?
\end{enumerate}
The corresponding calculated observables are the static and quasienergy dispersions, their bandwidths, site probabilities $P_j(t)$, mean position $\avg{x}(t)$, root-mean-square width $\sigma_x(t)$ and excited-band probability $P_{\rm ex}(mT)$. Bandwidth is a spectral quantity; width and center are real-space observables; momentum distributions provide an experimental sign-sensitive connection.

\section{Historical context}
Bloch theory expresses spatial translation symmetry through momentum-labeled wave functions. Quantum formulations of time-periodic dynamics by Shirley and Sambe established the analogous quasienergy and extended-space methods \cite{Shirley1965,Sambe1973}. Dunlap and Kenkre derived dynamic localization for an ac-driven tight-binding lattice in 1986 \cite{Dunlap1986}; Holthaus connected strong ac forcing to miniband collapse in 1992 \cite{Holthaus1992}. Madison et al. reported optical-lattice Bloch-band suppression in 1998 \cite{Madison1998}.

The tunability of ultracold gases made these ideas accessible to matter-wave transport. Lignier et al. measured controllable tunneling suppression in 2007 \cite{Lignier2007}. Eckardt et al. studied sinusoidal and square-wave forcing in 2009 \cite{Eckardt2009}. Holthaus's later tutorial provides the broader framework for computing continuous-lattice quasienergy bands, including effects beyond an isolated band \cite{Holthaus2016}. This report follows that progression: spatial bands, time-periodic bands, transport, and limits of the approximation.

\chapter{Stationary optical lattice and Bloch bands}
\section{From a standing wave to the Hamiltonian}
Let the complex electric-field amplitude be
\begin{equation}
 \mathcal E(x)=\mathcal E_b e^{\ii k_{\mathrm L}x}+\mathcal E_b e^{-\ii k_{\mathrm L}x}
 =2\mathcal E_b\cos(k_{\mathrm L}x),
\end{equation}
where the physical field is $\operatorname{Re}[\mathcal E(x)e^{-\ii\omega_{\mathrm L}t}]$ and $\omega_{\mathrm L}$ is the optical angular frequency. In the scalar far-detuned approximation, a real dynamic polarizability $\alpha_{\rm pol}$ gives
\begin{equation}
 U_{\rm dip}(x)=-\frac14\alpha_{\rm pol}|\mathcal E(x)|^2
 =-\alpha_{\rm pol}|\mathcal E_b|^2\cos^2(k_{\mathrm L}x).
\end{equation}
For red detuning, choose $V_0=\alpha_{\rm pol}|\mathcal E_b|^2>0$. Adding the irrelevant constant $V_0$ gives $V_0\sin^2(k_{\mathrm L}x)$. Blue-detuned fields can directly give positive-intensity barriers after a shift of the spatial origin. Spontaneous scattering and vector or tensor light shifts are neglected. The motional Hamiltonian is therefore
\begin{equation}\label{eq:H0}
 H_0=\frac{p^2}{2m}+V_0\sin^2(k_{\mathrm L}x),\qquad p=-\ii\hbar\partial_x.
\end{equation}
Both terms have units of energy. The one-dimensional stationary Schr\"odinger equation is $H_0\phi=E\phi$.

\section{Reciprocal lattice and Bloch states}
The notation in this section and the next follows Jos\'e Men\'endez's PHY 581 lecture notes, particularly Lecture 7\_1, Eqs.~(19)--(32), and Lecture 8\_2, Eqs.~(1)--(2). In those notes, $\bm q$ labels a general plane wave, $\bm k$ is chosen in the first Brillouin zone, and $\bm K,\bm K'$ and $\bm G$ are reciprocal lattice vectors. The potential is $U(\bm r)$, the full eigenfunction is $\psi_{n\bm k}(\bm r)$, and its energy is $\mathcal E_{n\bm k}$. For our one-dimensional lattice these vectors reduce to scalar wave numbers: $\bm r\to x$, $\bm k\to k$, $\bm K\to K$ and $\bm G\to G$. The sample volume and primitive-cell volume reduce to $L$ and $d$, respectively. Here $U(x)=V(x)=V_0\sin^2(k_{\mathrm L}x)$ and $\mathcal H=H_0$.

\paragraph{Reciprocal lattice and allowed plane waves.}
Since $U(x+d)=U(x)$ with $d=\pi/k_{\mathrm L}$, a reciprocal wave number $K$ satisfies
\begin{equation}
 e^{\ii Kd}=1,\qquad K_\ell=\frac{2\pi\ell}{d}=2\ell k_{\mathrm L},\qquad \ell\in\mathbb Z.
\end{equation}
Thus $\ell$ is simply an integer label for a reciprocal lattice vector. For a sample of $N_c$ cells, $L=N_cd$, impose $\psi(x+L)=\psi(x)$. A plane wave then obeys $e^{\ii qL}=1$, which gives $q=2\pi r/L$ for integer $r$. Following the notes, its normalized basis function is
\begin{equation}\label{eq:notesbasis}
 \langle x|q\rangle\equiv\phi_q(x)=\frac{1}{\sqrt L}e^{\ii qx},\qquad
 \langle q|q'\rangle=\frac1L\int_0^L e^{\ii(q'-q)x}\dd x=\delta_{qq'}.
\end{equation}
In these two sections $q$ denotes a plane-wave wave number, with units of inverse length, as in the notes. The dimensionless momentum $q=kd$ used elsewhere in the report and in the code is explicitly identified again at the end of the derivation.

Every allowed plane-wave wave number can be written
\begin{equation}\label{eq:notesdecomposition}
 q=k-K,\qquad -k_{\mathrm L}\leq k<k_{\mathrm L},
\end{equation}
with a reciprocal lattice wave number $K$. This merely folds $q$ into one chosen Brillouin zone. For example, a wave with $q=k+2k_{\mathrm L}$ corresponds to $K=-2k_{\mathrm L}$. At the zone boundary one uses the stated half-open interval to avoid counting the same reduced momentum twice.

\paragraph{Bloch expansion in the lecture notation.}
The periodic potential couples plane waves only within one fixed-$k$ block, as derived in Section 2.3. Its energy eigenfunctions can therefore be expanded as
\begin{equation}\label{eq:pw}
 \psi_{nk}(x)=\sum_K c^{(n)}_{k-K}e^{\ii(k-K)x}
 =e^{\ii kx}\underbrace{\sum_K c^{(n)}_{k-K}e^{-\ii Kx}}_{u_{nk}(x)}.
\end{equation}
Here $c^{(n)}_{k-K}$ is the Fourier coefficient multiplying the plane wave of wave number $k-K$. With this convention the factor $L^{-1/2}$ is absorbed into $c^{(n)}_{k-K}$, just as in the lecture's wave-function expansion. For a wave function normalized over the sample, orthogonality gives
\begin{equation}
 \int_0^L|\psi_{nk}(x)|^2\dd x=L\sum_K|c^{(n)}_{k-K}|^2=1.
\end{equation}
Equivalently, the coefficients of the normalized kets $|k-K\rangle$ are $\sqrt L\,c^{(n)}_{k-K}$ and have unit sum of squared magnitudes.

For every direct-lattice translation $R=jd$ with integer $j$, $e^{-\ii KR}=1$. Consequently,
\begin{align}
 u_{nk}(x+R)&=\sum_K c^{(n)}_{k-K}e^{-\ii Kx}e^{-\ii KR}=u_{nk}(x),\\
 \psi_{nk}(x+R)&=e^{\ii kR}\psi_{nk}(x).
\end{align}
This establishes the Bloch form in the same notation as the notes:
\begin{equation}\label{eq:bloch}
 \boxed{\psi_{nk}(x)=e^{\ii kx}u_{nk}(x),\qquad u_{nk}(x+d)=u_{nk}(x),\qquad
 \mathcal H\psi_{nk}=\mathcal E_{nk}\psi_{nk}.}
\end{equation}
At one fixed $k$, diagonalization produces multiple eigenvalues $\mathcal E_{nk}$ and their eigenfunctions. The band index $n$ distinguishes those eigenstates; the reciprocal index $K$ distinguishes the plane-wave components inside each eigenstate. The exact expansion contains infinitely many reciprocal vectors.

For consistency with the rest of the report, $\psi_{nk}$ here is the state denoted by $\phi_{nk}$ elsewhere, and $\mathcal E_{nk}=E_n(k)$. Uppercase $K$ in these two sections is a reciprocal lattice wave number with units of inverse length. The later dimensionless shaking strength $K=F_0d/(\hbar\omega)$ is a separate, explicitly defined quantity.

\section{Deriving the matrix eigenvalue problem}
\paragraph{Kinetic energy in the plane-wave basis.}
Write the stationary Hamiltonian in the lecture notation as
\begin{equation}
 \mathcal H=\hat T+U(x),\qquad \hat T=-\frac{\hbar^2}{2m}\frac{\dd^2}{\dd x^2}.
\end{equation}
The first and second derivatives of a normalized basis function are
\begin{equation}
 \frac{\dd\phi_{q'}}{\dd x}=\ii q'\phi_{q'},\qquad
 \frac{\dd^2\phi_{q'}}{\dd x^2}=-(q')^2\phi_{q'}.
\end{equation}
Therefore $\hat T\phi_{q'}=\hbar^2(q')^2\phi_{q'}/(2m)$ and
\begin{align}
 \langle q|\hat T|q'\rangle
 &=\frac{\hbar^2(q')^2}{2m}\frac1L\int_0^L e^{\ii(q'-q)x}\dd x\nonumber\\
 &=\frac{\hbar^2q^2}{2m}\delta_{qq'}.
\end{align}
In the last line $q'=q$ whenever the Kronecker delta is nonzero. Replacing $q=k-K$ and $q'=k-K'$ gives
\begin{equation}\label{eq:noteskinetic}
 \boxed{\langle k-K|\hat T|k-K'\rangle=
 \frac{\hbar^2(k-K)^2}{2m}\delta_{KK'}.}
\end{equation}
Thus the kinetic energy is diagonal in this basis: differentiation multiplies a plane wave by a number and preserves its wave number.

\paragraph{Potential Fourier coefficients and the selection rule.}
As in the notes, expand the periodic potential using reciprocal lattice vectors:
\begin{equation}\label{eq:notespotential}
 U(x)=\sum_G U_G e^{\ii Gx},\qquad
 U_G=\frac1d\int_0^d U(x)e^{-\ii Gx}\dd x.
\end{equation}
The Fourier coefficient $U_G$ has units of energy. Between normalized plane waves,
\begin{align}
 \langle q|U|q'\rangle
 &=\frac1L\int_0^L e^{-\ii qx}\left[\sum_GU_Ge^{\ii Gx}\right]e^{\ii q'x}\dd x\nonumber\\
 &=\sum_GU_G\frac1L\int_0^L e^{\ii(G+q'-q)x}\dd x\nonumber\\
 &=\sum_GU_G\delta_{q,G+q'}.
\end{align}
The integral vanishes unless $q-q'=G$, so a periodic potential can directly couple plane waves only when their wave numbers differ by a reciprocal lattice vector with a nonzero Fourier coefficient. In particular, two different reduced momenta $k$ and $k'$ in the half-open first Brillouin zone are not coupled. This is why the full matrix separates into independent blocks labeled by $k$.

Within one block, substitute $q=k-K$ and $q'=k-K'$:
\begin{align}
 \langle k-K|U|k-K'\rangle
 &=\sum_GU_G\delta_{k-K,G+k-K'}\nonumber\\
 &=\sum_GU_G\delta_{G,K'-K}=U_{K'-K}.
\end{align}
Combining this with \cref{eq:noteskinetic} gives precisely the lecture's matrix convention:
\begin{equation}\label{eq:notesmatrix}
 \boxed{H_{KK'}(k)\equiv\langle k-K|\mathcal H|k-K'\rangle
 =\frac{\hbar^2(k-K)^2}{2m}\delta_{KK'}+U_{K'-K}.}
\end{equation}
The row is labeled by $K$ and the column by $K'$. Since a real potential obeys $U_{-G}=U_G^*$, this matrix is Hermitian.

\paragraph{Collecting the coefficients of one basis wave.}
To see directly how the matrix acts on the expansion, substitute \cref{eq:pw} into $\mathcal H\psi_{nk}=\mathcal E_{nk}\psi_{nk}$. The kinetic term is
\begin{equation}
 \hat T\psi_{nk}=\sum_{K'}\frac{\hbar^2(k-K')^2}{2m}c^{(n)}_{k-K'}e^{\ii(k-K')x},
\end{equation}
while the potential term is
\begin{equation}
 U\psi_{nk}=\sum_{G,K'}U_Gc^{(n)}_{k-K'}e^{\ii(k-K'+G)x}.
\end{equation}
Fix one target wave $e^{\ii(k-K)x}$. In the kinetic term its coefficient comes from $K'=K$. In the potential term it comes from every pair satisfying $k-K'+G=k-K$, or $G=K'-K$. Equating the coefficient of this target wave on the two sides therefore yields
\begin{equation}\label{eq:notescoefficient}
 \boxed{\frac{\hbar^2(k-K)^2}{2m}c^{(n)}_{k-K}
 +\sum_{K'}U_{K'-K}c^{(n)}_{k-K'}=\mathcal E_{nk}c^{(n)}_{k-K}.}
\end{equation}
This equation holds for every reciprocal vector $K$. It is the eigenvalue problem $\sum_{K'}H_{KK'}(k)c^{(n)}_{k-K'}=\mathcal E_{nk}c^{(n)}_{k-K}$.

\paragraph{Specialization to the optical cosine lattice.}
For the potential of this project,
\begin{equation}
 U(x)=V_0\sin^2(k_{\mathrm L}x)=\frac{V_0}{2}
 -\frac{V_0}{4}e^{2\ii k_{\mathrm L}x}-\frac{V_0}{4}e^{-2\ii k_{\mathrm L}x}.
\end{equation}
Comparison with \cref{eq:notespotential} identifies all its nonzero Fourier coefficients:
\begin{equation}
 U_0=\frac{V_0}{2},\qquad U_{2k_{\mathrm L}}=U_{-2k_{\mathrm L}}=-\frac{V_0}{4}.
\end{equation}
All other $U_G$ vanish. A choice $U_0=0$ would subtract $V_0/2$ from every energy; here the mean is retained to preserve the report's energy reference and numerical tables.

Label reciprocal vectors by $K_\ell=2\ell k_{\mathrm L}$ and define the normalized basis amplitudes $a_{n\ell}(k)=\sqrt L\,c^{(n)}_{k-K_\ell}$. For a fixed row $\ell$, the only contributing columns are $K'=K_\ell$ and $K'=K_\ell\pm2k_{\mathrm L}$. Thus \cref{eq:notescoefficient} becomes
\begin{equation}\label{eq:pwmat}
 \boxed{\left[\frac{\hbar^2(k-2\ell k_{\mathrm L})^2}{2m}+\frac{V_0}{2}\right]a_{n\ell}
 -\frac{V_0}{4}(a_{n,\ell-1}+a_{n,\ell+1})=\mathcal E_{nk}a_{n\ell}.}
\end{equation}
Define $D_\ell=\hbar^2(k-2\ell k_{\mathrm L})^2/(2m)+V_0/2$ and $g=V_0/4$.
The full matrix, ordered by increasing $\ell$, is
\begin{equation}
 H(k)=\begin{pmatrix}
 \ddots&\ddots&&&&&\\
 \ddots&D_{-2}&-g&&&&\\
 &-g&D_{-1}&-g&&&\\
 &&-g&D_0&-g&&\\
 &&&-g&D_1&-g&\\
 &&&&-g&D_2&\ddots\\
 &&&&&\ddots&\ddots
 \end{pmatrix},\qquad
 H(k)\begin{pmatrix}\vdots\\a_{n,-2}\\a_{n,-1}\\a_{n,0}\\a_{n,1}\\a_{n,2}\\\vdots\end{pmatrix}
 =\mathcal E_{nk}\begin{pmatrix}\vdots\\a_{n,-2}\\a_{n,-1}\\a_{n,0}\\a_{n,1}\\a_{n,2}\\\vdots\end{pmatrix}.
\end{equation}
Every blank entry is zero. ``Tridiagonal'' describes its three nonzero diagonals; the exact matrix still has infinitely many rows. A $3\times3$ example keeps only $\ell=-1,0,1$, whereas the numerical calculation uses $\ell=-M,\ldots,M$ with $M=12$, giving 25 rows. This finite cutoff omits the couplings beyond its endpoints and must be checked by increasing $M$. The index $\ell$ labels momentum components, while $n$ labels different eigenvectors of the complete matrix at a fixed $k$.

Divide \cref{eq:pwmat} by $\ER=\hbar^2k_{\mathrm L}^2/(2m)$ and define $s=V_0/\ER$ and $\kappa=k/k_{\mathrm L}$. The dimensionless matrix equation is
\begin{equation}\label{eq:staticdimensionless}
 \sum_{\ell'}h^{(0)}_{\ell\ell'}(\kappa)a_{n\ell'}=\frac{\mathcal E_{nk}}{\ER}a_{n\ell},\qquad
 h^{(0)}_{\ell\ell'}=\left[(\kappa-2\ell)^2+\frac{s}{2}\right]\delta_{\ell\ell'}
 -\frac{s}{4}(\delta_{\ell',\ell+1}+\delta_{\ell',\ell-1}).
\end{equation}
The code's basis $e^{\ii(k+2\ell_{\rm code}k_{\mathrm L})x}$ is obtained by relabeling $\ell_{\rm code}=-\ell$. Equivalently its coefficient $c_{n,\ell_{\rm code}}$ equals $a_{n,-\ell_{\rm code}}$. This reversal permutes rows and columns and preserves every eigenvalue. The code's dimensionless momentum is $q=kd=\pi\kappa$; it differs from the general plane-wave wave number $q$ used earlier in this derivation.

\paragraph{The first gap at the zone edge.}
At $s=0$, the diagonal entries give the folded free-particle energies $\hbar^2(k-K)^2/(2m)$. At $k=k_{\mathrm L}$, the waves with $K=0$ and $K=2k_{\mathrm L}$ have wave numbers $+k_{\mathrm L}$ and $-k_{\mathrm L}$ and the same kinetic energy $\ER$. Keeping this degenerate pair at weak lattice depth gives
\begin{equation}
 H_{\rm edge}=\begin{pmatrix}\ER+U_0&U_{2k_{\mathrm L}}\\U_{-2k_{\mathrm L}}&\ER+U_0\end{pmatrix}
 =\begin{pmatrix}\ER+V_0/2&-V_0/4\\-V_0/4&\ER+V_0/2\end{pmatrix}.
\end{equation}
Its characteristic equation and the two energies are
\begin{equation}
 (\ER+V_0/2-\mathcal E)^2-(V_0/4)^2=0,\qquad
 \mathcal E_\pm=\ER+\frac{V_0}{2}\pm\frac{V_0}{4}.
\end{equation}
The leading weak-lattice gap is therefore $\mathcal E_+-\mathcal E_-=V_0/2$. Additional plane waves change the gap at larger depth; the full finite-basis calculation below includes their mixing.

\fig{01_static_bands}{Static Bloch bands from the plane-wave matrix. The horizontal coordinate is $k/k_{\mathrm L}$, covering one Brillouin zone; the vertical coordinate is $E_n/\ER$. The lowest three bands are shown for depths $s=2,4,8,12$, using $M=12$ and 401 momentum points. The lowest band is the solid teal curve. Its bandwidth decreases with increasing depth, while the first interband gap grows. Dashed and dotted curves are the next two bands. These are numerical continuum results.}{fig:static}

\section{Calculated static quantities and a known benchmark}
Define the lowest-band width and edge gap by
\begin{equation}
 W_0=\max_k E_0(k)-\min_k E_0(k),\qquad \Delta_{01}^{\rm edge}=E_1(k_{\mathrm L})-E_0(k_{\mathrm L}).
\end{equation}
The following numbers are produced by the scripts, not transcribed as experimental values:
\begin{table}[!htbp]\centering\small
\begin{tabular}{rrrrrr}\toprule
$s$&$E_0(0)/\ER$&$E_0(k_{\mathrm L})/\ER$&$W_0/\ER$&$\Delta_{01}^{\rm edge}/\ER$&$J_1/\ER$\\\midrule
2&0.878234&1.470654&0.592420&0.996112&0.142755\\
4&1.544861&1.889751&0.344890&1.969357&0.085490\\
8&2.486043&2.609323&0.123280&3.769876&0.030799\\
12&3.165608&3.214620&0.049012&5.304419&0.012252\\\bottomrule
\end{tabular}
\caption{Reproduced lowest-band parameters. The definition of $J_1$ is derived below. Digits shown exceed what is needed to interpret the plots and allow comparison with the saved arrays.}
\label{tab:static}\end{table}
For an independent check, set $z=k_{\mathrm L}x$ in the stationary equation. It becomes
\begin{equation}
 \phi''(z)+\left[\frac{E}{\ER}-\frac{s}{2}+\frac{s}{2}\cos(2z)\right]\phi(z)=0.
\end{equation}
After $z\mapsto z+\pi/2$ this is the standard Mathieu equation with characteristic parameter $s/4$. Its lowest edges are
\begin{equation}
 \frac{E_0(0)}{\ER}=\frac{s}{2}+a_0(s/4),\qquad
 \frac{E_0(k_{\mathrm L})}{\ER}=\frac{s}{2}+b_1(s/4),
\end{equation}
where $a_0$ and $b_1$ are Mathieu characteristic values. Agreement with their independent SciPy evaluation is checked in the code. The widths and edge gaps at $s=4,8$ also reproduce Holthaus's rounded band parameters \cite{Holthaus2016}.

\section{Wannier states and the tight-binding approximation}
For one isolated band, choose normalized Wannier states
\begin{equation}
 \ket{w_j}=\frac{1}{\sqrt{N_c}}\sum_k e^{-\ii kjd}\ket{\phi_{0k}},\qquad
 \ket{\phi_{0k}}=\frac{1}{\sqrt{N_c}}\sum_j e^{\ii kjd}\ket{w_j}.
\end{equation}
They are orbitals centered at lattice cells, obtained by coherently combining the extended Bloch states. Smooth choices of Bloch phases give localized Wannier states. Translation invariance makes their Hamiltonian matrix elements depend only on separation. For a real inversion-symmetric lattice, define
\begin{equation}
 \bar E=\bra{w_j}H_0\ket{w_j},\qquad J_\ell=-\bra{w_{j+\ell}}H_0\ket{w_j}\in\mathbb R,\quad \ell\geq1.
\end{equation}
The band-projected Hamiltonian is
\begin{equation}\label{eq:longrangeH}
 H_{\rm band}=\bar E\sum_j\ket{w_j}\bra{w_j}
 -\sum_{\ell=1}^{\infty}J_\ell\sum_j(\ket{w_{j+\ell}}\bra{w_j}+\ket{w_j}\bra{w_{j+\ell}}).
\end{equation}
Acting on the inverse Wannier transform gives a factor $e^{-\ii\ell q}$ for a shift to $j+\ell$ and $e^{\ii\ell q}$ for its conjugate. Consequently,
\begin{align}
 E_0(q)&=\bar E-\sum_{\ell\geq1}J_\ell(e^{\ii\ell q}+e^{-\ii\ell q})\nonumber\\
 &=\bar E-2\sum_{\ell\geq1}J_\ell\cos(\ell q).\label{eq:fourierband}
\end{align}
Orthogonality of cosines gives a direct numerical extraction:
\begin{equation}\label{eq:hoppingextract}
 \bar E=\frac{1}{2\pi}\int_{-\pi}^{\pi}E_0(q)\dd q,\qquad
 J_\ell=-\frac{1}{2\pi}\int_{-\pi}^{\pi}E_0(q)\cos(\ell q)\dd q.
\end{equation}
Keeping only $J_1\equiv J>0$ gives the nearest-neighbor result
\begin{equation}\label{eq:cosine}
 H_{\rm NN}=\bar E\,I-J\sum_j(\ket{j+1}\bra j+\ket j\bra{j+1}),\qquad E_{\rm NN}(q)=\bar E-2J\cos q,
\end{equation}
whose bandwidth is $4J$. Here $\ket j$ abbreviates $\ket{w_j}$.

This approximation requires an isolated, initially occupied band and sufficiently localized Wannier functions that $|J_{\ell>1}|\ll J_1$. It is usually better in deeper lattices. Near a potential minimum, $\sin^2(k_{\mathrm L}x)\simeq k_{\mathrm L}^2x^2$; equating $V_0 k_{\mathrm L}^2x^2$ with $m\omega_{\rm well}^2x^2/2$ gives $\hbar\omega_{\rm well}=2\sqrt{s}\ER$. This estimates the well-level separation, while tunneling depends on overlap between different wells. The harmonic approximation alone does not calculate $J$.

\fig{02_tight_binding}{Nearest-neighbor accuracy for $s=2$ and $s=8$. The numerical lowest band (solid) and $-2\cos q$ (dashed) are plotted after subtracting $\bar E$ and dividing by the numerically extracted $J_1$. The horizontal coordinate is $k/k_{\mathrm L}$. Agreement improves with depth. The largest absolute energy errors divided by the full static bandwidth are $9.22\%$ and $1.07\%$, respectively; these are dispersion errors, not experimental uncertainties.}{fig:tb}

At $s=8$, the extracted next-neighbor hopping is $J_2/\ER=-6.35096\times10^{-4}$, about $2.06\%$ of $J_1$ in magnitude. A small correction in the static band becomes important when driving suppresses the leading term almost to zero. Therefore a deep lattice supports an accurate cosine approximation away from collapse without making every prediction at collapse exact.

\chapter{Periodically shaken lattices and Floquet theory}
\section{Laboratory Hamiltonian and moving coordinates}
Let the optical lattice displacement be $x_0(t)=X\cos(\omega t)$, where $X$ has units of length. A change in the relative optical phase produces this displacement. The laboratory Hamiltonian is
\begin{equation}\label{eq:lab}
 H_{\rm lab}(t)=\frac{p^2}{2m}+V_0\sin^2[k_{\mathrm L}(x-x_0(t))].
\end{equation}
Introduce $y=x-x_0(t)$ and the translated wave function $\chi(y,t)=\psi_{\rm lab}(y+x_0(t),t)$. The chain rule gives
\begin{align}
 \ii\hbar\partial_t\chi
 &=\ii\hbar[\partial_t\psi_{\rm lab}+\dot x_0\partial_x\psi_{\rm lab}]\nonumber\\
 &=\left[\frac{p_y^2}{2m}+V(y)-\dot x_0p_y\right]\chi,\qquad p_y=-\ii\hbar\partial_y.\label{eq:translated}
\end{align}
The momentum term appears because the coordinates move. To remove it, perform a velocity boost and scalar phase transformation:
\begin{equation}\label{eq:boost}
 \chi(y,t)=\exp\!\left(\frac{\ii m\dot x_0 y}{\hbar}\right)
 \exp\!\left[\frac{\ii m}{2\hbar}\int_0^t\dot x_0(t')^2\dd t'\right]\psi_F(y,t).
\end{equation}
For $B=\exp(\ii m\dot x_0y/\hbar)$,
\begin{equation}
 B^\dagger p_y B=p_y+m\dot x_0,\qquad
 -\ii\hbar B^\dagger\dot B=m\ddot x_0y.
\end{equation}
The kinetic and translation terms then give
\begin{equation}
 \frac{(p_y+m\dot x_0)^2}{2m}-\dot x_0(p_y+m\dot x_0)
 =\frac{p_y^2}{2m}-\frac{m\dot x_0^2}{2}.
\end{equation}
The derivative of the second phase in \cref{eq:boost} contributes $+m\dot x_0^2/2$, canceling the scalar term. Therefore
\begin{equation}\label{eq:forceH}
 \ii\hbar\partial_t\psi_F=H_F(t)\psi_F,\qquad
 H_F(t)=\frac{p_y^2}{2m}+V(y)-F(t)y,
\end{equation}
with inertial force
\begin{equation}\label{eq:force}
 F(t)=-m\ddot x_0(t)=mX\omega^2\cos(\omega t)\equiv F_0\cos(\omega t).
\end{equation}
The force points opposite to the lattice acceleration. The force amplitude has units $\mathrm{kg\,m\,s^{-2}}$, and $F_0d$ is an energy.

\section{A gauge that restores explicit spatial periodicity}
The linear potential in $H_F$ obscures ordinary spatial periodicity. Define the momentum-valued function
\begin{equation}
 A(t)=\int_0^t F(t')\dd t'=\frac{F_0}{\omega}\sin(\omega t)=-m\dot x_0(t),\qquad \dot A=F.
\end{equation}
Set $\psi_F=e^{\ii A(t)y/\hbar}\psi_A$. Since
\begin{equation}
 e^{-\ii Ay/\hbar}p_y e^{\ii Ay/\hbar}=p_y+A,\qquad
 -\ii\hbar e^{-\ii Ay/\hbar}\partial_t e^{\ii Ay/\hbar}=\dot A\,y,
\end{equation}
the two terms proportional to $y$ cancel. The resulting Hamiltonian is
\begin{equation}\label{eq:vectorH}
 H_A(t)=\frac{[p_y+A(t)]^2}{2m}+V(y).
\end{equation}
It is periodic in both space and time. The symbol $A$ is a momentum shift for a neutral-atom problem, not an assumed electromagnetic vector potential acting on a charge.

It is useful to keep track of the scalar phase convention. Direct translation in \cref{eq:translated} gives $H_0-\dot x_0p_y=H_A-m\dot x_0^2/2$. Their Floquet spectra differ by the mean quiver energy $mX^2\omega^2/4=\ER K^2/(2\pi^2)$. This is independent of $k$, so it does not affect bandwidths, derivatives or probabilities. All continuum computations below retain the $A^2/(2m)$ term in $H_A$.

\section{Floquet's theorem and quasienergy equivalence}
For a finite basis representation of a $T$-periodic Hermitian Hamiltonian, the one-period propagator is unitary. Let
\begin{equation}
 U(T,0)\ket{v_\alpha}=e^{-\ii\varepsilon_\alpha T/\hbar}\ket{v_\alpha}.
\end{equation}
Define
\begin{equation}
 \ket{u_\alpha(t)}=e^{\ii\varepsilon_\alpha t/\hbar}U(t,0)\ket{v_\alpha}.
\end{equation}
Periodicity of $H$ implies $U(t+T,0)=U(t,0)U(T,0)$, so $u_\alpha(t+T)=u_\alpha(t)$. The physical solution is
\begin{equation}\label{eq:floquet}
 \ket{\psi_\alpha(t)}=e^{-\ii\varepsilon_\alpha t/\hbar}\ket{u_\alpha(t)}.
\end{equation}
Substitute this into the time-dependent Schr\"odinger equation:
\begin{align}
 \ii\hbar\partial_t[e^{-\ii\varepsilon t/\hbar}u]
 &=e^{-\ii\varepsilon t/\hbar}(\varepsilon u+\ii\hbar\partial_tu)=e^{-\ii\varepsilon t/\hbar}Hu,\\
 (H-\ii\hbar\partial_t)u&=\varepsilon u.\label{eq:Sambe}
\end{align}
For any integer $r$, the replacements $\varepsilon\mapsto\varepsilon+r\hbar\omega$ and $u(t)\mapsto e^{\ii r\omega t}u(t)$ leave $\psi(t)$ unchanged. Thus quasienergy is an eigenphase expressed as an energy, defined modulo $\hbar\omega$. It is not generally the cycle-averaged expectation of $H$, nor a quantity to be minimized as an ordinary equilibrium ground-state energy. A reduced quasienergy zone can be chosen as $[-\hbar\omega/2,\hbar\omega/2)$.

\section{Floquet--Bloch states}
At each $k$, write
\begin{equation}
 \psi_{\alpha k}(y,t)=e^{\ii ky}e^{-\ii\varepsilon_\alpha(k)t/\hbar}u_{\alpha k}(y,t),\quad
 u_{\alpha k}(y+d,t)=u_{\alpha k}(y,t+T)=u_{\alpha k}(y,t).
\end{equation}
The cell-periodic Hamiltonian is
\begin{equation}\label{eq:Hk}
 H_k(t)=\frac{[p_y+\hbar k+A(t)]^2}{2m}+V(y).
\end{equation}
Spatial Bloch theory factors out $e^{\ii kx}$ and diagonalizes $H_k$ on one cell. Floquet theory factors out $e^{-\ii\varepsilon t/\hbar}$ and diagonalizes $H_k-\ii\hbar\partial_t$ on one cell and one period. The band label $\alpha$ identifies a Floquet family, while copies shifted by $r\hbar\omega$ represent the same physical state. In a multiband spectrum, ordering reduced quasienergies by size does not reliably track a single family through avoided crossings.

\section{Equations for a numerical Floquet calculation}
Define
\begin{equation}
 \tau_R=\ER t/\hbar,\quad \Omega_R=\hbar\omega/\ER,\quad
 K=\frac{F_0d}{\hbar\omega},\quad a=\frac{A_{\max}}{\hbar k_{\mathrm L}}=\frac{K}{\pi}.
\end{equation}
The symbol $a$ in this section is a dimensionless momentum amplitude, not the lattice spacing. In the plane-wave basis, the time-dependent matrix is
\begin{equation}\label{eq:continuumtime}
 h_{\ell\ell'}(q,\tau_R)=\left[(\kappa+2\ell+a\sin\Omega_R\tau_R)^2+\frac{s}{2}\right]\delta_{\ell\ell'}-\frac{s}{4}(\delta_{\ell',\ell+1}+\delta_{\ell',\ell-1}),
\end{equation}
and its evolution equation is $\ii\partial_{\tau_R}\bm c=h(q,\tau_R)\bm c$. The dimensionless period is $\mathcal T_R=2\pi/\Omega_R$.

One numerical method evolves each basis vector through one period to obtain $U_q(\mathcal T_R,0)$, then diagonalizes it. If $\lambda_\alpha$ is one of its eigenvalues,
\begin{equation}\label{eq:eigenphase}
 \frac{\varepsilon_\alpha(q)}{\ER}=-\frac{\Omega_R}{2\pi}\operatorname{Arg}\lambda_\alpha.
\end{equation}
This is the propagator approach emphasized by Holthaus \cite{Holthaus2016}. We use the exponential midpoint approximation
\begin{equation}\label{eq:midpoint}
 U_q\simeq\prod_{r=N_t-1}^{0}\exp[-\ii h(q,(r+1/2)\Delta\tau_R)\Delta\tau_R],\quad \Delta\tau_R=\mathcal T_R/N_t.
\end{equation}
The latest-time factor is on the left. Each exponential is unitary. Global time error is second order for a fixed finite matrix, but it must still be checked by increasing $N_t$.

An alternative is a Fourier expansion in time, $u=\sum_{\ell,n}C_{\ell n}e^{2\ii\ell k_{\mathrm L}y}e^{\ii n\omega t}$. If $h(t)=\sum_r h^{(r)}e^{\ii r\omega t}$, substitution into \cref{eq:Sambe} gives
\begin{equation}\label{eq:SambeMatrix}
 \sum_{\ell',m}\left[h^{(n-m)}_{\ell\ell'}+n\Omega_R\delta_{\ell\ell'}\delta_{nm}\right]C_{\ell'm}
 =\frac{\varepsilon}{\ER}C_{\ell n}.
\end{equation}
To obtain every nonzero Fourier coefficient, put $P_\ell=\kappa+2\ell$ and expand
\begin{equation}
 (P_\ell+a\sin\theta)^2=P_\ell^2+\frac{a^2}{2}+2aP_\ell\sin\theta-\frac{a^2}{2}\cos2\theta.
\end{equation}
It follows that
\begin{align}
 h^{(0)}_{\ell\ell'}&=\left(P_\ell^2+s/2+a^2/2\right)\delta_{\ell\ell'}-\frac{s}{4}(\delta_{\ell',\ell+1}+\delta_{\ell',\ell-1}),\\
 h^{(+1)}_{\ell\ell'}&=-\ii aP_\ell\delta_{\ell\ell'},\qquad h^{(-1)}_{\ell\ell'}=+\ii aP_\ell\delta_{\ell\ell'},\\
 h^{(+2)}_{\ell\ell'}&=h^{(-2)}_{\ell\ell'}=-a^2\delta_{\ell\ell'}/4.
\end{align}
All other coefficients vanish. A Fourier calculation must converge both the spatial cutoff $M$ and the temporal cutoff $|n|\leq N_F$. The eigenproblem \cref{eq:SambeMatrix} is provided to reproduce or extend the theory; the delivered multiband numerical results use \cref{eq:midpoint}, not an independently executed Sambe calculation.

\chapter{Effective tunneling and control of transport}
\section{The driven nearest-neighbor equation}
Project the inertial-force Hamiltonian onto one Wannier band and approximate the position by its cell center: $y\ket j\simeq jd\ket j$. This discards interband position matrix elements and intrawell structure. With $\bar E$ removed as a common phase, the site amplitudes $c_j(t)$ obey
\begin{equation}\label{eq:forceTB}
 \ii\hbar\dot c_j=-J(c_{j+1}+c_{j-1})-F(t)jd\,c_j.
\end{equation}
Define the dimensionless impulse
\begin{equation}
 \chi(t)=\frac{d}{\hbar}\int_0^tF(t')\dd t'=K\sin\omega t,\qquad K=\frac{F_0d}{\hbar\omega}=\frac{mX\omega d}{\hbar}.
\end{equation}
For comparison with an experimental displacement,
\begin{equation}
 \frac{X}{d}=\frac{2K}{\pi^2\Omega_R},\qquad \Omega_R=\frac{\hbar\omega}{\ER}.
\end{equation}
Take $c_j=e^{\ii j\chi(t)}b_j$. The time derivative and neighbor terms are
\begin{align}
 \ii\hbar\dot c_j&=e^{\ii j\chi}(\ii\hbar\dot b_j-\hbar j\dot\chi\,b_j),\\
 c_{j\pm1}&=e^{\ii j\chi}e^{\pm\ii\chi}b_{j\pm1}.
\end{align}
Since $\hbar\dot\chi=F(t)d$, substitution into \cref{eq:forceTB} cancels the terms proportional to $j$. Thus
\begin{equation}\label{eq:Peierls}
 \ii\hbar\dot b_j=-J(e^{\ii\chi}b_{j+1}+e^{-\ii\chi}b_{j-1}).
\end{equation}
The phase change is diagonal in position, so $|b_j|^2=|c_j|^2$ at all times. The phase of a hopping amplitude has become periodic.

Use the unitary Fourier transform on the infinite lattice,
\begin{equation}
 b_j(t)=\frac{1}{\sqrt{2\pi}}\int_{-\pi}^{\pi}b(q,t)e^{\ii qj}\dd q,\qquad
 b(q,t)=\frac{1}{\sqrt{2\pi}}\sum_jb_j(t)e^{-\ii qj}.
\end{equation}
Shifting $j$ by $\pm1$ multiplies the Fourier component by $e^{\pm\ii q}$. Therefore
\begin{equation}\label{eq:dispersiontime}
 \ii\hbar\partial_tb(q,t)=-2J\cos[q+K\sin(\omega t)]\,b(q,t),
\end{equation}
and direct integration gives
\begin{equation}\label{eq:exactq}
 b(q,t)=b(q,0)\exp\left[\frac{2\ii J}{\hbar}\int_0^t\cos(q+K\sin\omega t')\dd t'\right].
\end{equation}
Every $q$ evolves independently; the quasimomentum weights $|b(q,t)|^2$ are constant in this gauge. In the force gauge, the instantaneous wave number is shifted by the impulse: $k(t)=k+A(t)/\hbar$, obeying $\hbar\dot k(t)=F(t)$.

\section{Deriving the Bessel renormalization}
The phase accumulated in one period defines a continuous quasienergy representative:
\begin{equation}\label{eq:quasiaverage}
 \varepsilon(q)-\bar E=-\frac{2J}{T}\int_0^T\cos[q+K\sin(\omega t)]\dd t.
\end{equation}
Set $\theta=\omega t$ and use $T\omega=2\pi$:
\begin{align}
 \frac{1}{T}\int_0^T\cos(q+K\sin\omega t)\dd t
 &=\frac{1}{2\pi}\int_0^{2\pi}\big[\cos q\cos(K\sin\theta)-\sin q\sin(K\sin\theta)\big]\dd\theta.\label{eq:besselintegral}
\end{align}
Under $\theta\mapsto\theta+\pi$, $\sin\theta$ changes sign, so the second integral cancels between the two half cycles. The integral representation of the ordinary Bessel function is
\begin{equation}
 \mathcal J_0(K)=\frac{1}{2\pi}\int_0^{2\pi}e^{\ii K\sin\theta}\dd\theta
 =\frac{1}{2\pi}\int_0^{2\pi}\cos(K\sin\theta)\dd\theta.
\end{equation}
Combining these expressions produces
\begin{equation}\label{eq:Besselresult}
 \boxed{\varepsilon(q)=\bar E-2J\mathcal J_0(K)\cos q,\qquad \Jeff=J\mathcal J_0(K),\qquad W=4|\Jeff|.}
\end{equation}
This result is exact for the homogeneous, noninteracting, nearest-neighbor lattice. It is not merely the first term of a high-frequency expansion in that model. Indeed, with a periodic or infinite lattice, the shift operators $S\ket j=\ket{j+1}$ and $S^\dagger$ commute. Hamiltonians of the form $-J(e^{-\ii\chi}S+e^{\ii\chi}S^\dagger)$ consequently commute at different times. Their time-ordered exponential equals the exponential of their time integral. Finite open boundaries break this argument at the edges; the delivered periodic lattice is chosen large enough that the packet never reaches them.

High driving frequency relative to tunneling becomes relevant when additional terms do not commute, as with interactions or confinement. Using the model for real atoms also requires small interband transfer. A useful initial scale separation is $J\ll\hbar\omega\ll\Delta_{01}$, but it is not a sufficient guarantee: the force amplitude and multiphoton resonances must also be checked. The exactly solvable model itself permits any nonzero $\omega$ at fixed $K$.

\section{Narrowing, collapse and inversion}
For weak forcing, expanding the Bessel function gives
\begin{equation}
 \mathcal J_0(K)=1-\frac{K^2}{4}+\frac{K^4}{64}+O(K^6).
\end{equation}
Thus $W\simeq4J(1-K^2/4)$ initially decreases. At the first positive zero,
\begin{equation}
 K_*=2.4048255577\ldots,\qquad \mathcal J_0(K_*)=0,
\end{equation}
the band is flat in the nearest-neighbor model. For $K_*<K<5.5200781103\ldots$, $\mathcal J_0(K)<0$. The curvature at $q=0$ changes sign and the minima of the chosen continuous dispersion move from $q=0$ to $q=\pm\pi$. This is band inversion in the present project; it does not denote a topological band inversion involving different orbital bands.

\fig{03_floquet_bands}{Nearest-neighbor quasienergy bands at $K=0,1.5,K_*$ and $3.5$. Horizontal axes show $k/k_{\mathrm L}$; vertical axes show $(\varepsilon-\bar E)/J$. Solid curves are the analytical Bessel result; circles are independent numerical period averages using 2048 time samples. They illustrate the original band, narrowing, collapse and inversion in that order. These smooth representatives can be placed without folding in a quasienergy zone of width $\hbar\omega=8J$. The formula is independent of frequency at fixed $K$ within this model.}{fig:floquet}

\fig{04_tunneling_bandwidth}{The signed tunneling $\Jeff/J=\mathcal J_0(K)$ (upper panel) and nonnegative normalized bandwidth $W/(4J)=|\mathcal J_0(K)|$ (lower panel). The horizontal axis is the dimensionless drive strength. Vertical dotted lines mark the first two Bessel zeros. A bandwidth measurement alone cannot determine whether the band has inverted. These are evaluated analytical predictions, without fitted parameters.}{fig:tunneling}

For a fixed small $k$, the effective mass is determined by curvature:
\begin{equation}
 \frac1{m^*(k)}=\frac1{\hbar^2}\frac{\partial^2\varepsilon}{\partial k^2}
 =\frac{2\Jeff d^2}{\hbar^2}\cos(kd).
\end{equation}
Near $k=0$, its sign follows $\Jeff$. At collapse, the inverse mass vanishes in this approximation. A uniform bipartite chain permits the transformation $b_j\mapsto(-1)^jb_j$, which flips the sign of a real hopping and shifts quasimomentum by $\pi$. The operational sign comparison must therefore hold the initial state and momentum reference fixed. The directed-motion calculation below does so; it does not re-prepare a new band minimum for each drive strength.

\section{Quasienergy slope and average velocity}
There are two complementary derivations. In the exact single-band model, the instantaneous velocity of a narrow packet is the static group velocity at its accelerated momentum:
\begin{equation}
 v_k(t)=\frac{1}{\hbar}\frac{\partial E_0(k+A/\hbar)}{\partial k}.
\end{equation}
Differentiating its period-averaged quasienergy gives
\begin{equation}
 \overline v_k=\frac1T\int_0^T v_k(t)\dd t
 =\frac1\hbar\frac{\partial}{\partial k}\left[\frac1T\int_0^TE_0(k+A/\hbar)\dd t\right]
 =\frac1\hbar\frac{\partial\varepsilon(k)}{\partial k}.\label{eq:velocity}
\end{equation}
The nearest-neighbor expression is
\begin{equation}\label{eq:velocityNN}
 \boxed{\overline v_k=\frac{2\Jeff d}{\hbar}\sin(kd).}
\end{equation}
This has units length/time because $\Jeff d/\hbar$ has those units. A packet at $k=0$ has zero mean velocity for either sign of $\Jeff$, so it cannot by itself reveal the inversion.

For a general isolated Floquet branch, use the extended-space inner product
\begin{equation}
 \langle\!\langle u|v\rangle\!\rangle=\frac1T\int_0^T\dd t\int_0^d\dd y\,u^*(y,t)v(y,t),\quad \langle\!\langle u|u\rangle\!\rangle=1.
\end{equation}
Differentiate $(H_k-\ii\hbar\partial_t)u_k=\varepsilon_k u_k$ with respect to $k$, take the inner product with $u_k$, and use Hermiticity and periodic temporal boundary conditions. The terms containing $\partial_k u_k$ cancel, leaving
\begin{align}
 \partial_k\varepsilon_k&=\langle\!\langle u_k|\partial_k H_k|u_k\rangle\!\rangle,\\
 \partial_kH_k&=\hbar\frac{p_y+\hbar k+A(t)}{m}=\hbar\hat v_y.
\end{align}
This proves \cref{eq:velocity} as a Floquet Hellmann--Feynman relation, for a normalized differentiable branch. Degeneracies require choosing smooth eigenstates appropriately. The mean laboratory velocity is the same because the added lattice velocity has zero cycle average. A packet occupying many $k$ values has the weighted mean of the branch velocities, rather than the slope at a single arbitrarily chosen $k$ \cite{Arlinghaus2011}.

\section{Relative phases and dynamic localization}
At integer periods $t=mT$, the periodic functions in \cref{eq:floquet} return to their initial values. In the nearest-neighbor model, $\chi(mT)=0$ and
\begin{equation}
 b(q,mT)=b(q,0)e^{-\ii[\varepsilon(q)-\bar E]mT/\hbar}.
\end{equation}
Two components $q$ and $q'$ acquire the relative phase
\begin{equation}
 \Delta\varphi_{q,q'}(mT)=-\frac{mT}{\hbar}[\varepsilon(q)-\varepsilon(q')].
\end{equation}
If the band is dispersive, this phase depends on the pair of momenta. Their recombination in real space changes with time, producing drift or spreading. At $K=K_*$, every component acquires the same phase. Hence
\begin{equation}\label{eq:revival}
 b_j(mT)=b_j(0),\qquad c_j(mT)=c_j(0),
\end{equation}
apart from the removed common phase $e^{-\ii\bar E mT/\hbar}$. Every initial single-band packet revives, not only the examples in the figures. This is dynamic localization in the ideal model \cite{Dunlap1986,Eckardt2009}.

The revival does not imply immobility at all intermediate times. To show this explicitly, define the dimensionless time and integrated hopping
\begin{equation}\label{eq:Bdef}
 \tau_J=Jt/\hbar,\quad \Omega_J=\hbar\omega/J,\quad
 B(\tau_J)=\int_0^{\tau_J}e^{\ii K\sin(\Omega_Ju)}\dd u.
\end{equation}
Taking the real part of $e^{\ii q}B$ in \cref{eq:exactq} gives
\begin{equation}\label{eq:phaseB}
 b(q,\tau_J)=b(q,0)\exp\{2\ii\operatorname{Re}[e^{\ii q}B(\tau_J)]\}.
\end{equation}
The Jacobi--Anger expansion, $e^{\ii K\sin\theta}=\sum_{n\in\mathbb Z}\mathcal J_n(K)e^{\ii n\theta}$, can be integrated term by term:
\begin{equation}\label{eq:Bseries}
 B(\tau_J)=\mathcal J_0(K)\tau_J+
 \sum_{n\ne0}\mathcal J_n(K)\frac{e^{\ii n\Omega_J\tau_J}-1}{\ii n\Omega_J}.
\end{equation}
The second term is periodic and vanishes at complete cycles. At $K_*$, only this bounded term remains. Wave-packet deformation within a period is the micromotion; its cancellation at successive periods is the localization. In laboratory coordinates the lattice displacement must also be added to a packet center, $\avg{x}_{\rm lab}=x_0(t)+\avg y$. This displacement is periodic and has no effect on the central width.

\section{Probability distribution and an exact width identity}
For an initial particle on site zero, $b_j(0)=\delta_{j0}$, the Fourier amplitude is $1/\sqrt{2\pi}$. Put $B=|B|e^{\ii\beta}$ in \cref{eq:phaseB}; the inverse transform gives
\begin{align}
 b_j(\tau_J)&=\frac{1}{2\pi}\int_{-\pi}^{\pi}e^{\ii qj}e^{2\ii|B|\cos(q+\beta)}\dd q\nonumber\\
 &=\ii^j e^{-\ii j\beta}\mathcal J_j(2|B|),\qquad
 P_j(\tau_J)=\mathcal J_j^2(2|B|).\label{eq:probBessel}
\end{align}
The second equality follows by expanding $e^{\ii z\cos\theta}=\sum_n\ii^n\mathcal J_n(z)e^{\ii n\theta}$ and retaining the Fourier term $n=-j$. At $B=0$, the limiting state is again $\delta_{j0}$.

The first moment vanishes because $P_j=P_{-j}$. To derive the second moment without assuming a continuum Gaussian, use $j\leftrightarrow\ii\partial_q$ and integrate by parts over the periodic zone. For the normalized Fourier wave function,
\begin{align}
 \sum_jj^2P_j&=\int_{-\pi}^{\pi}|\partial_qb(q)|^2\dd q,\\
 \partial_qb(q)&=-2\ii|B|\sin(q+\beta)b(q),\\
 \sum_jj^2P_j&=\frac{4|B|^2}{2\pi}\int_{-\pi}^{\pi}\sin^2(q+\beta)\dd q=2|B|^2.
\end{align}
Thus
\begin{equation}\label{eq:widthdelta}
 \boxed{\sigma_x(t)/d=\sqrt2|B(\tau_J)|,\qquad
 \sigma_x(mT)/d=\sqrt2\,|\mathcal J_0(K)|\,\tau_J(mT).}
\end{equation}
At zero forcing this is ballistic spreading; at collapse it is a bounded periodic width with zero value at every strobe. The finite width of a real Wannier orbital would add intrawell structure to the density. The plotted site width concerns the probability of occupying different lattice cells.

\section{A broad Gaussian packet and the sign of tunneling}
For a packet with position standard deviation $\sigma_0\gg d$ and central wave number $k_0$, take
\begin{equation}\label{eq:gaussian}
 b_j(0)=\mathcal N\exp\!\left[-\frac{(jd)^2}{4\sigma_0^2}+\ii k_0jd\right],\qquad \sum_j|b_j(0)|^2=1,
\end{equation}
with normalization $\mathcal N$. The factor 4 ensures that the probability variance is approximately $\sigma_0^2$. The quasimomentum probability variance is approximately $1/(4\sigma_0^2)$ in physical $k$ units. Expanding a smooth quasienergy near $k_0$,
\begin{equation}
 \varepsilon(k_0+\delta k)=\varepsilon_0+\varepsilon_0'\delta k+\tfrac12\varepsilon_0''(\delta k)^2+\cdots,
\end{equation}
the linear term translates the envelope by $\varepsilon_0't/\hbar$, while the quadratic phase gives
\begin{equation}\label{eq:Gaussianwidth}
 \avg{x}(t)\simeq\frac{\varepsilon_0'}{\hbar}t,\qquad
 \sigma_x^2(t)\simeq\sigma_0^2+\left(\frac{\varepsilon_0''t}{2\hbar\sigma_0}\right)^2.
\end{equation}
For example, in the momentum representation $\hat x(t)=\ii\partial_k+t\varepsilon'(k)/\hbar$. With an initially real Gaussian momentum envelope, the symmetrized position--velocity covariance is zero. Linearizing the velocity and using $\operatorname{Var}(\delta k)=1/(4\sigma_0^2)$ gives \cref{eq:Gaussianwidth} directly. This is a narrow-momentum, smooth-band approximation; it does not replace the exact lattice propagation used for the figures.

For a finite-width packet the exact stroboscopic center in the nearest-neighbor model obeys
\begin{equation}\label{eq:packetvelocity}
 \frac{\avg{x}(mT)-\avg{x}(0)}{d}=2\mathcal J_0(K)\tau_J(mT)\int_{-\pi}^{\pi}|b(q,0)|^2\sin q\,\dd q.
\end{equation}
For the broad Gaussian this average is approximately $e^{-d^2/(8\sigma_0^2)}\sin(k_0d)$, obtained from the characteristic function of a Gaussian momentum distribution. The code evaluates the discrete weighted average directly. A packet with $k_0d=\pi/3$ then moves in opposite directions when $\Jeff$ changes sign, while its initial state is held fixed.

\section{Longer-range hopping and incomplete collapse}
Retaining \cref{eq:fourierband} changes the instantaneous dispersion to
\begin{equation}
 E_0[q+K\sin\omega t]=\bar E-2\sum_{\ell\geq1}J_\ell\cos[\ell q+\ell K\sin\omega t].
\end{equation}
Applying the same integral used in \cref{eq:besselintegral} to each harmonic gives
\begin{equation}\label{eq:longrangequasi}
 \boxed{\varepsilon_{\rm iso}(q)=\bar E-2\sum_{\ell\geq1}J_\ell\mathcal J_0(\ell K)\cos(\ell q).}
\end{equation}
This is exact for the homogeneous single-band model with all its hopping ranges. As a model of the continuous optical lattice, it assumes negligible interband transitions and neglects any nontrivial geometric phase along the momentum path. For the present inversion-symmetric one-dimensional cosine lattice the sinusoidal path retraces itself, so the single-band geometric phase cancels.

At $K=K_*$ the $\ell=1$ term vanishes, but generally $\mathcal J_0(\ell K_*)\ne0$ for $\ell>1$. A nonzero $J_2$ therefore leaves dispersion and slow residual spreading. Band collapse in the nearest-neighbor approximation does not establish exact collapse for a cosine optical potential. This distinction is part of the established analysis of Eckardt et al. \cite{Eckardt2009}.

Square-wave forcing provides a useful contrast connected to their experiment. If $F=F_0$ for half a period and $F=-F_0$ for the other half, choose a zero-mean impulse $\chi$ that sweeps linearly between $-\pi K/2$ and $+\pi K/2$. Both passages spend equal time in each impulse interval, so
\begin{align}
 \overline{e^{\ii\ell\chi}}&=\frac{1}{\pi K}\int_{-\pi K/2}^{\pi K/2}e^{\ii\ell\chi}\dd\chi
 =\frac{\sin(\ell\pi K/2)}{\ell\pi K/2},\\
 J_{\ell,\rm eff}&=J_\ell\operatorname{sinc}(\ell\pi K/2),\qquad \operatorname{sinc}z=\sin z/z.
\end{align}
At nonzero $K=2r$ with integer $r$, every nonzero integer harmonic vanishes, giving an exactly flat dispersion within the isolated-band model. This does not remove interband or interaction effects. No square-wave simulation is claimed here; the relation is derived for the experimental comparison.

\chapter{Numerical calculations and physical interpretation}
\section{Reproducible protocol and parameter choices}
All plotted numbers are generated by the delivered Python scripts. No figure reproduces an experimental trace. The static continuum calculation uses the potential \cref{eq:convention}; the principal transport calculations solve the specified nearest-neighbor model. The longer-range calculation uses numerically extracted static-band Fourier coefficients. Only the final comparison explicitly allows transfer between bands.

\begin{table}[!htbp]\centering\small
\begin{tabular}{P{0.36\linewidth}P{0.56\linewidth}}\toprule
Calculation&Parameters and numerical representation\\\midrule
Static bands&$s=2,4,8,12$; $M=12$; 401 momenta for plots; 2048 momenta for hopping integrals\\
Nearest-neighbor bands&$K=0,1.5,K_*,3.5$; numerical period average with 2048 samples\\
Wave packets&$N_s=1024$ periodic sites; $\Omega_J=4$; $0\leq t\leq20T$; 1601 output times; $|n|\leq40$ in \cref{eq:Bseries}\\
Directed packet&$\sigma_0=6d$; $k_0d=\pi/3$; same initial amplitudes for every $K$\\
Longer-range band&$s=2,8$; 20 hopping ranges; 2048 momenta for band widths\\
Multiband comparison&$s=8$; $K=0.35$; $k=0$; $M=8$; 2048 midpoint steps per period; $m=0,\ldots,20$\\\bottomrule
\end{tabular}\caption{Numerical settings. Continuum time uses $\tau_R=\ER t/\hbar$; nearest-neighbor time uses $\tau_J=Jt/\hbar$. They satisfy $\tau_J=(J/\ER)\tau_R$ and $\Omega_J=\Omega_R/(J/\ER)$.}\label{tab:protocol}\end{table}

The choices $K=1.5$ and $3.5$ give $\mathcal J_0(K)=0.51182767$ and $-0.38012774$. They were selected to show substantial narrowing and a clearly negative hopping. They are not optimized experimental settings. Because the single-band theory is independent of frequency at fixed $K$, $\Omega_J=4$ is chosen to make intra-cycle motion visible. These transport plots alone do not certify that a particular real optical lattice will remain in one band at all these settings.

\section{Fourier propagation and observables}
For $N_s$ sites with unit-normalized amplitudes, the discrete transforms used are
\begin{equation}
 b(q_r)=\frac1{\sqrt{N_s}}\sum_j b_j e^{-\ii q_rj},\qquad
 b_j=\frac1{\sqrt{N_s}}\sum_r b(q_r)e^{\ii q_rj},\quad q_r=\frac{2\pi r}{N_s}\pmod{2\pi}.
\end{equation}
NumPy's \texttt{norm="ortho"} implements this normalization. Centered site arrays are shifted before the forward transform and shifted back after the inverse transform. At each output time the code multiplies every momentum amplitude by the phase in \cref{eq:phaseB}. Thus it integrates the full driven dispersion rather than using $\Jeff$ to construct a continuously stationary approximation.

The observables are calculated directly from probabilities:
\begin{equation}\label{eq:moments}
 P_j=|b_j|^2,\quad \mathcal P=\sum_jP_j,\quad
 \frac{\avg{x}}d=\sum_jjP_j,\quad \frac{\sigma_x^2}{d^2}=\sum_jj^2P_j-\left(\sum_jjP_j\right)^2.
\end{equation}
The norm $\mathcal P$ is checked rather than renormalized during evolution. Because the physical wave function is approximately $\sum_jc_j(t)w_j(x)$, these are site populations. For well-localized Wannier orbitals the continuous density is approximately $\sum_jP_j|w_j(x)|^2$ if overlap terms are negligible; resolving interference or a time-of-flight pattern requires coherent sums of amplitudes, not this population approximation.

\section{Spreading and complete-cycle revival}
\fig{05_probability_evolution}{Site-resolved wave-packet evolution from $b_j(0)=\delta_{j0}$. Each horizontal axis is the site coordinate $j=x/d$ and each vertical axis is elapsed time in periods. The four drive strengths are the same as in \cref{fig:floquet}; $\Omega_J=4$. All panels use the identical linear probability scale $0\leq P_j\leq1$. The undriven and detuned packets spread, while the packet at $K_*$ stays near its initial site and revives after each period. The very low-probability wings are intentionally less visible on this full scale; \cref{fig:wings} displays them explicitly. The curves come from numerical Fourier propagation of the full time-dependent nearest-neighbor model.}{fig:probability}

\fig{05b_expanding_wings}{The same probability data for $K=0$ and $K=3.5$, with a linear color scale saturating at $P_j=0.05$. Axes and times are as in \cref{fig:probability}; values above the stated threshold share the darkest color. This view reveals the ballistic outward fronts and the slower spreading of the inverted band. Symmetric expansion does not reveal the negative sign of $\Jeff$ because it depends on its magnitude. The distribution for $K=1.5$ is intermediate in spreading speed.}{fig:wings}

\fig{06_width_micromotion}{RMS site width versus time for the localized initial state. In the left panel, continuous lines are numerical widths for $K=0,1.5,K_*,3.5$, and open markers at integer periods are the exact analytical stroboscopic prediction $\sqrt2|\mathcal J_0(K)|\tau_J$. The right panel resolves two periods at $K_*$: the width is nonzero within the cycle and returns to zero at $t/T=0,1,2$. Both vertical axes are $\sigma_x/d$; both horizontal axes are $t/T$. With $\Omega_J=4$, the maximum intra-cycle width is about $0.833d$.}{fig:width}

\fig{07_localization_snapshots}{Probabilities on individual sites within the first period at $K_*$ and $\Omega_J=4$. The horizontal coordinate is integer site $j$ and the vertical coordinate is $P_j$. The five curves correspond to $t/T=0,1/4,1/2,3/4,1$. The initial and one-period distributions coincide at a single occupied site; the intermediate distributions have neighboring-site population. Overlap between some curves reflects symmetry or revival, not missing calculations. The figure is generated by the same full driven propagation as the probability maps.}{fig:snapshots}

For zero drive, the exact nearest-neighbor result is $P_j=\mathcal J_j^2(2\tau_J)$, which gives ballistic fronts and $\sigma_x/d=\sqrt2\tau_J$. At $K=1.5$ the stroboscopic expansion rate is reduced by $0.51182767$. At $K=3.5$ it is reduced by $0.38012774$, even though the quasienergy is inverted. At collapse the calculated maximum width during a period is $0.832754d$, while the amplitude error of complete-cycle revivals over 20 periods is below $4\times10^{-15}$ in this ideal-model computation. That roundoff-level agreement verifies the implementation against the exact identity; it does not set the accuracy of a real optical-lattice experiment.

\section{A sign-sensitive directed-motion experiment in the model}
\fig{08_directed_motion}{Directed motion for an initially broad Gaussian with $\sigma_0=6d$ and $k_0d=\pi/3$. In the left panel the vertical coordinate is the mean position in units of $d$ and the horizontal coordinate is time in periods. Continuous lines include micromotion; open markers give the weighted group-velocity prediction at complete cycles. The $K=3.5$ packet moves toward negative $x$, whereas the undriven and $K=1.5$ packets move toward positive $x$. In the right panel, the distributions $P_j(20T)$ for $K=0$ and $3.5$ show their centers on opposite sides. The same initial state and $\Omega_J=4$ are used in every run.}{fig:drift}

The exact discrete initial momentum average is $\avg{\sin q}=0.86302359$. Predicted and fitted stroboscopic velocities, in units $dJ/\hbar$, are
\begin{equation}
 \begin{array}{c|rrrr}
 K&0&1.5&K_*&3.5\\\hline
 \overline v/(dJ/\hbar)&1.726047&0.883439&0&-0.656118.
 \end{array}
\end{equation}
The fits use centers at $m=0,\ldots,20$ and agree with \cref{eq:packetvelocity} to below $2\times10^{-14}$ sites in the tested trajectories. At $k_0=0$ these signed velocities would all vanish. Holding a nonzero initial momentum fixed is therefore essential to this demonstration.

\section{How much collapse survives longer-range hopping?}
\fig{09_long_range_collapse}{Residual quasienergy bandwidth near the first nearest-neighbor collapse point. Horizontal axes show $K$; vertical axes show $W(K)/W(0)$. Solid curves retain 20 numerically extracted hopping ranges in the isolated-band expression \cref{eq:longrangequasi}; dashed curves retain only the nearest neighbor and give $|\mathcal J_0(K)|$. The dotted vertical line marks $K_*$. Depths $s=2$ and $s=8$ are shown in separate panels; the inset uses the same axes quantities and enlarges the small residual at $s=8$. The residual width is smaller for the deeper lattice but remains nonzero. These are isolated-band calculations and do not include interband transitions.}{fig:longrange}

For each $K$, the code evaluates \cref{eq:longrangequasi} on 2048 momenta and takes the difference between its maximum and minimum; it does not assume that all extrema stay at the original zone center and edge. At $K_*$ it gives
\begin{equation}
 \begin{array}{c|rr}
 s&W(K_*)/\ER&W(K_*)/W(0)\\\hline
 2&0.0231503&0.0390775\\
 8&0.000616239&0.00499868.
 \end{array}
\end{equation}
The $s=8$ band retains about $0.50\%$ of its original width, whereas the shallow $s=2$ band retains about $3.91\%$. These residuals exceed the tested Fourier reconstruction error by many orders of magnitude. They reflect omitted hopping physics rather than numerical failure. A small residual dispersion eventually produces spreading over sufficiently long observation times; nearest-neighbor revivals cannot then be asserted indefinitely.

\section{Bounded multiband comparison}
The previous correction still assumes that the particle stays in one Bloch band. To test a different failure mechanism, propagate the continuous Hamiltonian \cref{eq:continuumtime} at $s=8$, $q=0$ and $K=0.35$. Start in the normalized static lowest-band eigenvector $\ket{g_0}$ at $t=0$, where $A(0)=0$. The drive amplitude is present from the start, with the stated sinusoidal phase; an adiabatic turn-on is not simulated. Compare
\begin{equation}
 \Omega_R=1\qquad\hbox{and}\qquad
 \Omega_R=\frac{E_1(0)-E_0(0)}{\ER}=5.18618959\ldots.
\end{equation}
The second drive is resonant with the center-of-zone first interband gap. At integer periods $A(mT)=0$, so the total probability outside the static lowest band is unambiguously
\begin{equation}\label{eq:Pexc}
 P_{\rm ex}(mT)=1-\left|\bra{g_0}U_q(T,0)^m\ket{g_0}\right|^2.
\end{equation}
At intermediate times an instantaneous accelerated-band projection would be needed to distinguish reversible basis motion from excitation. We restrict this diagnostic to complete cycles.

\fig{10_multiband_excitation}{Multiband excitation at $s=8$, $K=0.35$ and $k=0$. The horizontal axes are completed periods $m$ and the vertical axes are total excited-band probability at $mT$. Solid markers are continuum propagation with 17 plane waves ($M=8$) and 2048 unitary midpoint steps per period. The dashed zero baseline is the isolated-band prediction. The left panel is off resonance at $\Omega_R=1$ and remains below $2.56\times10^{-4}$. The right panel is at the static interband resonance $\Omega_R=5.18619$ and reaches about $0.997$. Its dotted curve is a resonant two-state rotating-wave approximation, derived below. The two vertical scales differ intentionally to resolve the small off-resonant excitation.}{fig:multiband}

Within the tested 20 periods, $\max_m P_{\rm ex}$ is $2.55822\times10^{-4}$ off resonance and $0.997106$ on resonance. At $20T$ the probabilities are $1.25073\times10^{-4}$ and $0.262337$, respectively. The large resonant excitation is coherent population transfer in a closed single-particle calculation. It is not a measured heating rate, irreversible thermalization or a prediction for a whole finite-width condensate. The calculation covers one quasimomentum and two frequencies; it does not establish a global safe-driving region.

For physical interpretation, expand the dimensionless continuum Hamiltonian in $a=K/\pi$:
\begin{equation}
 h(\tau_R)=h_0+2a\hat P\sin(\Omega_R\tau_R)+a^2\sin^2(\Omega_R\tau_R)I,\qquad \hat P=p/(\hbar k_{\mathrm L}).
\end{equation}
The scalar term can be removed as a common phase. In the static eigenbasis, the off-diagonal element between the two lowest states is $2aP_{10}\sin(\Omega_R\tau_R)$, where $|P_{10}|=0.96518139$ in the calculation. Moving to the interaction picture multiplies it by $e^{\ii(E_1-E_0)\tau_R/\ER}$. At resonance,
\begin{equation}
 2aP_{10}\sin(\Omega_R\tau_R)e^{\ii\Omega_R\tau_R}
 =\frac{aP_{10}}{\ii}(e^{2\ii\Omega_R\tau_R}-1).
\end{equation}
Discard the rapidly oscillating term and retain the constant coupling of magnitude $a|P_{10}|$. A two-state Hamiltonian with off-diagonal magnitude $g=a|P_{10}|$ has propagator $\cos(g\tau_R)I-\ii\sin(g\tau_R)H/g$, giving
\begin{equation}\label{eq:RWA}
 P_{\rm ex}^{\rm RWA}(\tau_R)=\sin^2(a|P_{10}|\tau_R).
\end{equation}
This rotating-wave approximation omits counterrotating terms, higher levels and drive-induced resonance shifts. It is an explanatory comparison, not the method used to generate the continuum data.

The numerical one-period propagator is also diagonalized using a complex Schur decomposition. The largest overlap weight of $\ket{g_0}$ with one Floquet eigenvector is about $0.999931$ off resonance and $0.526520$ on resonance. The latter is consistent with strong hybridization of two Floquet families. Reduced eigenphases and weights are saved in \texttt{results/multiband\_summary.json}; they are not mislabeled as a globally tracked lowest quasienergy band.

\section{Convergence checks and known limits}
The code verifies the mathematical implementation by several independent comparisons. The summarized errors below are maximum absolute differences over the stated finite samples, not rigorous bounds on an infinite-dimensional problem.
\begin{table}[!htbp]\centering\small
\begin{tabular}{P{0.59\linewidth}r}\toprule
Check and tested domain&Observed difference\\\midrule
Lowest three static bands; $M=6$ vs $12$, four depths, 33 momenta&$<7\times10^{-12}\ER$\\
Mathieu lowest-band edges; four depths&$<3\times10^{-13}\ER$\\
Free-particle limit $s=0$; five bands, 33 momenta&zero to recorded precision\\
20-range Fourier reconstruction; $s=2$, 401 momenta&$<7\times10^{-10}\ER$\\
Period quadrature vs Bessel formula; sampled $q,K$&$<3\times10^{-16}J$\\
Integrated-hopping series vs adaptive quadrature&$<3\times10^{-15}$ in $\tau_J$\\
Real-space ODE vs FFT amplitudes; $K=1.5$, $t=3T$&$<6\times10^{-15}$\\
Undriven continuum propagator vs static spectral evolution&$<2\times10^{-13}$\\
Continuum $P_{\rm ex}$; 2048 vs 4096 steps, off resonance&$<9\times10^{-9}$\\
Continuum $P_{\rm ex}$; 2048 vs 4096 steps, on resonance&$<8\times10^{-7}$\\
Continuum $P_{\rm ex}$; $M=8$ vs $10$, both frequencies&$<2\times10^{-11}$\\\bottomrule
\end{tabular}
\caption{Numerical checks in the completed run. Propagator errors use the matrix spectral norm. All excitation-probability refinements compare the full discrete trace $m=0,\ldots,20$. Saved JSON gives full domains and unrounded errors.}\label{tab:checks}\end{table}

For static reconstruction the momentum integral is compared at 512, 1024 and 2048 momenta and at 10 and 20 hopping ranges. For packet propagation, 256, 512 and 1024 sites give the same tested unforced width to roundoff; a boundary probability below $3\times10^{-30}$ at 1024 sites supports absence of wraparound in the plotted delta-packet runs. The free Bessel distribution, probability conservation, zero mean position for the initial single-site packet, and exact complete-cycle revivals are separately checked. The real-space ODE uses a different representation from the Fourier propagator.

The continuum propagators remain unitary within about $5.1\times10^{-13}$ in spectral norm. Unitarity alone would not detect phase errors, so the time-step refinement is essential. Comparing 512, 1024 and 2048 steps to 4096 is consistent with the expected second-order discretization error before roundoff dominates. All assertions in \texttt{convergence.py} passed. These tests support the displayed finite calculations; they do not extend the model to interactions, different resonances or indefinitely long times.

\chapter{Connections to experiments}
\section{What an experiment can measure}
In situ absorption imaging gives a spatial density profile, from which cloud centers and widths can be extracted. Expansion along a lattice after releasing longitudinal confinement probes the transport generated by the dispersion. Momentum distributions after releasing the atoms provide complementary information about relative site phases and quasimomentum occupation. A phase-coherent lattice state produces interference peaks whose positions and contrast depend on those phases. These observables need not reconstruct every quasienergy directly to test the underlying band engineering.

For a narrow Gaussian envelope, \cref{eq:Gaussianwidth} shows that the asymptotic expansion rate depends on the magnitude of the band curvature. For a localized site state, \cref{eq:widthdelta} gives a rate proportional to $|\Jeff|$. Inferring a hopping magnitude from expansion therefore depends on the prepared packet and observation regime. It does not by itself measure the sign. The directed-motion demonstration here is a proposed model observable; the published experiments used their own preparation and detection protocols.

\section{Lignier et al. 2007}
Lignier and collaborators loaded a Bose--Einstein condensate into a one-dimensional optical lattice and shook it through a controlled frequency difference between the lattice beams \cite{Lignier2007}. Releasing longitudinal confinement allowed them to extract an effective tunneling magnitude from cloud expansion. The normalized result followed the predicted $|\mathcal J_0(K)|$ dependence, including strong suppression near its zeros. Their reported experimental resolution established suppression by at least a factor of 25, rather than a mathematically exact zero.

They also used time-of-flight interference to examine coherence and the negative-hopping regime. The observed momentum-pattern shift supported a sign change of the effective tunneling. Near collapse, reduced interwell coupling allowed faster interaction-driven dephasing. At large forcing frequencies, excited-band population was observed. Thus the paper supplies evidence for tunable hopping, its sign-sensitive consequences and limitations of a purely isolated-band description; it does not justify identifying every strongly shaken condensate with an ideal noninteracting chain.

\section{Eckardt et al. 2009}
Eckardt and collaborators compared sinusoidal and square-wave inertial forcing of a dilute condensate in a one-dimensional optical lattice \cite{Eckardt2009}. Their in situ expansion measurements tested the normalized effective hopping magnitude. The nearest-neighbor predictions were $|\mathcal J_0(K)|$ for sinusoidal forcing and $|\operatorname{sinc}(\pi K/2)|$ for square-wave forcing. They also measured the decay of time-of-flight interference visibility to locate rapid dephasing near band collapse.

The first observed suppression points were consistent with $K\simeq2.4$ for sinusoidal forcing and $K\simeq2.0$ for square-wave forcing. Their analysis distinguished incomplete sinusoidal collapse caused by longer-range hopping from exact square-wave collapse within an isolated-band description. The measured lattices were deep enough that longer-range hopping was below the stated experimental sensitivity. Our residual-bandwidth calculation illustrates that theoretical distinction; it is not a fit to their data or a new experimental test.

\section{What the present models reproduce}
\begin{table}[!htbp]\centering\small
\begin{tabular}{P{0.27\linewidth}P{0.31\linewidth}P{0.32\linewidth}}\toprule
Observable or effect&Established experimental connection&Calculation in this project\\\midrule
Expansion suppression&Lignier 2007; Eckardt 2009&Exact-model site probabilities and RMS widths\\
Signed tunneling&Momentum interference in Lignier 2007&Fixed-$k_0$ directed drift; no simulated experimental images\\
Drive-waveform dependence&Eckardt 2009, sinusoidal and square-wave forcing&Derived Bessel and sinc laws; sinusoidal numerical plots\\
Coherence loss near collapse&Interference visibility in both papers&Discussed; no interacting dephasing simulation\\
Interband excitation&Reported among limitations in Lignier 2007&Bounded coherent resonance example; no quantitative experimental fit\\\bottomrule
\end{tabular}\caption{Relation between published observations and the scope of the reproduced models. The entries identify mechanisms and observables; they do not assert numerical agreement between our parameters and experimental conditions.}\label{tab:experiment}\end{table}

\section{Additional modeling needed for quantitative comparison}
A quantitative comparison must first calibrate the actual lattice depth, spacing, shaking amplitude, frequency, drive phase and turn-on envelope. The force and displacement relations derived above convert experimental beam modulation into $K$; uncertainties in these parameters propagate into an inferred $\Jeff$. One must also model the initial cloud's momentum distribution rather than assume a site delta state or the particular Gaussian used here.

Weakly interacting condensate dynamics can be treated with a driven Gross--Pitaevskii equation,
\begin{equation}
 \ii\hbar\partial_t\Psi(y,t)=\left[\frac{p^2}{2m}+V(y)-F(t)y+V_{\rm trap}(y)+g_{1\rm D}|\Psi(y,t)|^2\right]\Psi(y,t).
\end{equation}
Here $\int|\Psi|^2\dd y=N_a$ is the atom number, $g_{1\rm D}$ has units energy times length, and $V_{\rm trap}$ is the residual longitudinal confinement. This equation is a mean-field model: it can describe interaction-induced density dynamics but does not fully capture number fluctuations or many-body phase diffusion.

For number fluctuations and strong correlations, an appropriate starting point is the driven Bose--Hubbard Hamiltonian \cite{Jaksch1998},
\begin{equation}
 H_{\rm BH}(t)=-J\sum_j(\hat a_{j+1}^\dagger\hat a_j+\mathrm{h.c.})+
 \frac{U}{2}\sum_j\hat n_j(\hat n_j-1)+\sum_j[V_j-F(t)jd]\hat n_j.
\end{equation}
The bosonic operator $\hat a_j$ removes an atom from site $j$, $\hat n_j=\hat a_j^\dagger\hat a_j$, $U$ is an onsite interaction energy and $V_j$ is the trap energy at that site. In an effective one-dimensional model $U=g_{1\rm D}\int|w_j(y)|^4\dd y$. Replacing $J$ by $\Jeff$ in this interacting Hamiltonian is a further effective approximation; the exact commuting-Hamiltonian argument used for a single particle no longer applies.

Higher-band occupation needs multiband propagation over the occupied range of $k$, with realistic ramps and densities. Finite temperature, atom loss, laser phase noise, spatial intensity envelopes and transverse modes may matter. Detection should also be modeled: in situ images convolve density with an imaging response, while ideal long time-of-flight images approximately measure
\begin{equation}
 n(p)\propto|\widetilde w(p)|^2\sum_{j,l}e^{\ii p(j-l)d/\hbar}\avg{\hat a_j^\dagger\hat a_l}.
\end{equation}
Here $\widetilde w(p)$ is the Wannier-orbital momentum envelope. The off-diagonal correlations determine interference contrast. Populations $P_j$ alone cannot predict the dephasing measurements. These additions are required before claiming quantitative agreement with a particular experimental trace.

\chapter{Conclusions}
\section{Answers to the central questions}
\textbf{How does shaking change a Bloch band?} For the homogeneous nearest-neighbor model, periodic forcing shifts the instantaneous quasimomentum and averages the dispersion over that excursion. The resulting quasienergy is $\bar E-2J\mathcal J_0(K)\cos(kd)$. Its bandwidth is $4J|\mathcal J_0(K)|$, its slope sets the cycle-averaged velocity, and its curvature sets the effective mass and narrow-packet spreading. This result follows from explicit phase integration and is exact within the stated model.

\textbf{When is spreading suppressed, and is the atom motionless?} At a Bessel zero the nearest-neighbor quasienergy is flat. All momentum components gain the same phase after a complete cycle, so the initial packet returns at every stroboscopic observation. The numerical distributions and exact width identity show reversible deformation within the cycle. Therefore dynamic localization concerns bounded motion and complete-cycle recurrence, not the absence of micromotion. A realistic lattice may retain slow spreading because its band is not a pure cosine.

\textbf{How is band inversion detected?} A negative $\mathcal J_0(K)$ reverses the sign of hopping and the slope at a fixed nonzero quasimomentum. Expansion from a symmetric localized packet measures a magnitude and cannot distinguish the sign. In the calculated fixed preparation with $k_0d=\pi/3$, the center velocity reverses from $+1.72605\,dJ/\hbar$ without driving to $-0.656118\,dJ/\hbar$ at $K=3.5$. Momentum interference provides the related published experimental connection. These statements hold with the momentum and preparation references fixed.

\textbf{Where does the simple description fail?} The cosine approximation becomes less accurate in shallow lattices and is particularly fragile when its dominant hopping is tuned to zero. The isolated-band calculation leaves residual bandwidths of about $3.91\%$ and $0.50\%$ at the first nominal collapse for $s=2$ and $8$. Separately, the bounded continuum propagation shows nearly complete transfer out of the initial lowest band near an interband resonance, even when the drive frequency greatly exceeds the tunneling scale. Interactions, confinement and noise introduce additional limitations that are absent from these ideal calculations.

\section{Physical significance and extensions}
The connection to solid-state physics is direct: spatial periodicity creates a band, and temporal periodicity creates a controllable quasienergy band. Shaking can regulate transport without changing the static barrier height. It can slow a packet, dynamically localize it within an ideal model or reverse its directed motion. The useful object is the entire dispersion and the prepared state's distribution over it, rather than an isolated parameter without its measurement protocol.

Lignier et al. and Eckardt et al. established matter-wave tests of the main suppression mechanism and its coherence-related consequences. The present project reproduces the corresponding ideal-model physics, exposes two distinct approximation limits and supplies all calculations in executable form. It presents neither a new localization mechanism nor a quantitative reproduction of those experimental datasets.

Natural extensions are to propagate a realistic finite-width packet in the multiband continuum, introduce a measured drive envelope, compare square-wave with sinusoidal forcing numerically, and include interactions and residual trapping. Systematic tracking of Floquet states through avoided crossings would connect the simple Bessel regime to more general band engineering. Any such extension needs renewed spatial, temporal and observation-time convergence checks; the present 20-period, one-momentum multiband comparison cannot settle indefinite-time stability of a continuous or interacting driven gas.

\appendix
\chapter{Symbols and dimensional checks}
{\small
\begin{longtable}{P{0.16\linewidth}P{0.56\linewidth}P{0.18\linewidth}}
\toprule Symbol&Meaning&Units\\\midrule\endfirsthead
\toprule Symbol&Meaning&Units\\\midrule\endhead
$x,y$&Laboratory and moving-frame position&length\\
$x_0,X$&Lattice displacement and its amplitude&length\\
$m$&Atomic mass; an integer strobe index is separately identified as $m$ in $mT$&mass\\
$k_{\mathrm L},\lambda_{\mathrm L}$&Optical wave number and wavelength; $k_{\mathrm L}=2\pi/\lambda_{\mathrm L}$&inverse length, length\\
$d$&Primitive lattice spacing $\pi/k_{\mathrm L}$&length\\
$G_\ell$&Reciprocal vector $2\ell k_{\mathrm L}$&inverse length\\
$k,\kappa,q$&Bloch wave number; $\kappa=k/k_{\mathrm L}$; $q=kd=\pi\kappa$&inverse length, 1, 1\\
$p,\hat P$&Momentum; dimensionless momentum $\hat P=p/(\hbar k_{\mathrm L})$&momentum, 1\\
$\ER$&Recoil energy $\hbar^2k_{\mathrm L}^2/(2m)$&energy\\
$V_0,s$&Lattice depth; scaled depth $s=V_0/\ER$&energy, 1\\
$E_n,\bar E$&Static band energy and lowest-band mean&energy\\
$J_\ell,J$&Hopping over $\ell$ cells; $J=J_1>0$&energy\\
$\Jeff$&Effective nearest-neighbor hopping $J\mathcal J_0(K)$&energy\\
$\omega,T$&Drive angular frequency; period $2\pi/\omega$&inverse time, time\\
$\omega_{\mathrm L}$&Optical carrier angular frequency&inverse time\\
$F,F_0$&Inertial force and its amplitude&force\\
$A,a$&Momentum impulse; $a=A_{\max}/(\hbar k_{\mathrm L})=K/\pi$&momentum, 1\\
$K,\chi$&Scaled drive strength $F_0d/(\hbar\omega)$; impulse phase $Ad/\hbar$&1\\
$\varepsilon_\alpha$&Floquet quasienergy, modulo $\hbar\omega$&energy\\
$\mathcal J_n$&Ordinary Bessel function of integer order $n$&1\\
$\tau_R,\tau_J$&Times $\ER t/\hbar$ and $Jt/\hbar$&1\\
$\Omega_R,\Omega_J$&Frequencies $\hbar\omega/\ER$ and $\hbar\omega/J$&1\\
$B$&Integrated hopping in \cref{eq:Bdef}&1\\
$P_j,\mathcal P$&Site probability and total probability&1\\
$\avg x,\sigma_x$&Mean position and standard deviation&length\\
$P_{\rm ex}$&Probability outside the initial static lowest band at a strobe&1\\
$N_c,N_s$&Cells for Bloch normalization; sites for packet FFT&integer count\\
$M,N_t,N_F$&Spatial plane-wave cutoff, time steps per period, temporal Fourier cutoff&integer count\\
\bottomrule
\end{longtable}
}
The symbol $m$ denotes atomic mass in Hamiltonians and an integer only when explicitly used as a strobe index. To check a central formula, $\partial_k\varepsilon$ has units energy times length, so $(1/\hbar)\partial_k\varepsilon$ is a velocity. Likewise $\varepsilon''t/(\hbar\sigma_0)$ has units length, and $F_0d/(\hbar\omega)$ is dimensionless. Standard deviations always refer to probability distributions, rather than unqualified Gaussian amplitude widths.

\chapter{Execution instructions and numerical provenance}
\section{Files and commands}
The complete commented Python files accompany this PDF in the project archive. They require Python 3.10 or later, NumPy, SciPy and Matplotlib; no external datasets or proprietary numerical tools are needed. The executed environment used Python 3.12.14, NumPy 2.3.5, SciPy 1.17.0 and Matplotlib 3.10.8. To reproduce all figures and checks from the extracted project folder:
\begin{lstlisting}
python -m venv .venv
source .venv/bin/activate
python -m pip install -r requirements.txt
python code/run_all.py
\end{lstlisting}
On Windows activate the environment with \texttt{.venv\textbackslash Scripts\textbackslash activate}. Individual files can also be run, for example \texttt{python code/static\_bands.py} or \texttt{python code/convergence.py}. The final driver stage raises an exception if a numerical acceptance condition fails. Running again overwrites generated outputs; the source files are not modified.

\begin{table}[!htbp]\centering\small
\begin{tabular}{P{0.35\linewidth}P{0.57\linewidth}}\toprule
Python file&Purpose\\\midrule
\texttt{common.py}&Units, tridiagonal static matrix, hopping extraction, shared plotting style\\
\texttt{static\_bands.py}&Static bands and parameters\\
\texttt{single\_band.py}&Bessel bands, hopping and bandwidth; longer-range collapse comparison\\
\texttt{wavepackets.py}&Full time-dependent Fourier propagation; distributions, width, micromotion, directed motion\\
\texttt{multiband.py}&Continuum one-period propagator, Floquet overlaps and excitation probabilities\\
\texttt{convergence.py}&Independent limits, norm checks, ODE comparison, spatial and temporal refinements\\
\texttt{run\_all.py}&Complete execution, software versions and SHA-256 output hashes\\\bottomrule
\end{tabular}\caption{Every numerical figure is associated with an executable source file. Shared definitions are imported from the same project, so paths do not depend on the shell's current directory.}\end{table}

Output subfolders contain vector PDF and PNG figures, compressed NumPy arrays, JSON summaries and a CSV static-band table. The exact numerical check domains and values are in \path{results/convergence.json}. \path{results/execution_record.json} records package versions, elapsed time, command arguments and hashes. A separate bounded execution in \texttt{validation/} records the input hashes before and after execution and its run status. These records document computation; they do not constitute a mathematical proof of all untested parameter regimes.

The PDF is typeset from \texttt{report.tex}. With a standard TeX installation, run the following command three times to resolve its contents and cross-references:
\begin{lstlisting}
pdflatex -interaction=nonstopmode -halt-on-error report.tex
\end{lstlisting}
No BibTeX run is needed because the bibliography is embedded. Python plotting does not require a LaTeX installation.

\section{Accuracy and coverage}
The static cutoff test covers four depths, 33 momentum values and the lowest three bands. It does not certify every high-energy band. The 20-hopping reconstruction is checked at 401 momenta for $s=2,8$. Packet tests cover the displayed drive strengths and times, supplemented by the zero-drive Bessel solution and one direct real-space ODE run. The continuum refinement covers only two frequencies at $s=8$, $K=0.35$, $k=0$ and the 21 stroboscopic observations. Temporal Fourier/Sambe truncation is derived but is not used for a second numerical Floquet spectrum.

No random sampling is used. Floating-point agreement at the $10^{-14}$ level is presented only where an independent identity or implementation supports it. Small nonmonotonic spatial-cutoff differences at the $10^{-12}$ level reflect numerical roundoff, rather than evidence of unresolved physical effects. Increasing the number of plane waves does not automatically establish uniform convergence of an unbounded continuous-lattice Floquet spectrum. Longer times and different resonance structures require separate checks.

\chapter{Source verification and attribution}
The principal published reference is Holthaus's 2016 tutorial \cite{Holthaus2016}. The primary texts actually checked for the detailed framework and experimental comparisons were arXiv:1510.09042v1, arXiv:0707.0403v1 and arXiv:0812.1997v2, together with publication metadata. Their dates differ from the publication years in the bibliography because the first and third are preprints or revised preprints. The accompanying \texttt{sources.md} records the exact versions, source links, claims used and notation translations. No licensed full paper is redistributed in the project archive.

The stationary and Floquet eigenproblems, Bessel and sinc renormalizations, group-velocity relation and dynamic-localization argument are established theory, derived explicitly here. The generated plots, finite parameter choices, error measurements and selected resonant example are reproduced computations. The discussion of requirements for experimental comparison is a modeling assessment. No digitized experimental points, fitted experimental parameters or novel theorem are asserted.

\begin{thebibliography}{99}
\addcontentsline{toc}{chapter}{References}
\bibitem{Holthaus2016} M. Holthaus, ``Floquet engineering with quasienergy bands of periodically driven optical lattices,'' \emph{J. Phys. B: At. Mol. Opt. Phys.} \textbf{49}, 013001 (2016). \href{https://doi.org/10.1088/0953-4075/49/1/013001}{doi:10.1088/0953-4075/49/1/013001}. Preprint: \href{https://arxiv.org/abs/1510.09042}{arXiv:1510.09042}.
\bibitem{Jaksch1998} D. Jaksch, C. Bruder, J. I. Cirac, C. W. Gardiner, and P. Zoller, ``Cold bosonic atoms in optical lattices,'' \emph{Phys. Rev. Lett.} \textbf{81}, 3108--3111 (1998). \href{https://doi.org/10.1103/PhysRevLett.81.3108}{doi:10.1103/PhysRevLett.81.3108}.
\bibitem{Bloch2008} I. Bloch, J. Dalibard, and W. Zwerger, ``Many-body physics with ultracold gases,'' \emph{Rev. Mod. Phys.} \textbf{80}, 885--964 (2008). \href{https://doi.org/10.1103/RevModPhys.80.885}{doi:10.1103/RevModPhys.80.885}.
\bibitem{Shirley1965} J. H. Shirley, ``Solution of the Schr\"odinger equation with a Hamiltonian periodic in time,'' \emph{Phys. Rev.} \textbf{138}, B979--B987 (1965). \href{https://doi.org/10.1103/PhysRev.138.B979}{doi:10.1103/PhysRev.138.B979}.
\bibitem{Sambe1973} H. Sambe, ``Steady states and quasienergies of a quantum-mechanical system in an oscillating field,'' \emph{Phys. Rev. A} \textbf{7}, 2203--2213 (1973). \href{https://doi.org/10.1103/PhysRevA.7.2203}{doi:10.1103/PhysRevA.7.2203}.
\bibitem{Dunlap1986} D. H. Dunlap and V. M. Kenkre, ``Dynamic localization of a charged particle moving under the influence of an electric field,'' \emph{Phys. Rev. B} \textbf{34}, 3625--3633 (1986). \href{https://doi.org/10.1103/PhysRevB.34.3625}{doi:10.1103/PhysRevB.34.3625}.
\bibitem{Holthaus1992} M. Holthaus, ``Collapse of minibands in far-infrared irradiated superlattices,'' \emph{Phys. Rev. Lett.} \textbf{69}, 351--354 (1992). \href{https://doi.org/10.1103/PhysRevLett.69.351}{doi:10.1103/PhysRevLett.69.351}.
\bibitem{Madison1998} K. W. Madison, M. C. Fischer, R. B. Diener, Q. Niu, and M. G. Raizen, ``Dynamical Bloch band suppression in an optical lattice,'' \emph{Phys. Rev. Lett.} \textbf{81}, 5093--5096 (1998). \href{https://doi.org/10.1103/PhysRevLett.81.5093}{doi:10.1103/PhysRevLett.81.5093}.
\bibitem{Lignier2007} H. Lignier, C. Sias, D. Ciampini, Y. Singh, A. Zenesini, O. Morsch, and E. Arimondo, ``Dynamical control of matter-wave tunneling in periodic potentials,'' \emph{Phys. Rev. Lett.} \textbf{99}, 220403 (2007). \href{https://doi.org/10.1103/PhysRevLett.99.220403}{doi:10.1103/PhysRevLett.99.220403}. Preprint: \href{https://arxiv.org/abs/0707.0403}{arXiv:0707.0403}.
\bibitem{Eckardt2009} A. Eckardt, M. Holthaus, H. Lignier, A. Zenesini, D. Ciampini, O. Morsch, and E. Arimondo, ``Exploring dynamic localization with a Bose--Einstein condensate,'' \emph{Phys. Rev. A} \textbf{79}, 013611 (2009). \href{https://doi.org/10.1103/PhysRevA.79.013611}{doi:10.1103/PhysRevA.79.013611}. Preprint: \href{https://arxiv.org/abs/0812.1997}{arXiv:0812.1997}.
\bibitem{Arlinghaus2011} S. Arlinghaus and M. Holthaus, ``Generalized acceleration theorem for spatiotemporal Bloch waves,'' \emph{Phys. Rev. B} \textbf{84}, 054301 (2011). \href{https://doi.org/10.1103/PhysRevB.84.054301}{doi:10.1103/PhysRevB.84.054301}.
\end{thebibliography}
\end{document}
